\paragraph{正规自同态的例：}

\begin{enumerate}
\def\labelenumi{\arabic{enumi})}
\item
  关于 \(\Phi\) 的酉自同构（\(\S 6\)，第 2 小节），它们由关系 \(u^{-1} = u^{*}\) 刻画（\(\S 1\)，第 8 小节，命题 8 的推论）；
\item
  自同态 \(u\)，满足 \(u^{*}=u\)；它们称为埃尔米特自同态。
\end{enumerate}

对 E 的每个埃尔米特自同态 \(u\)，令 \(\Phi_{u}(x,y)=\Phi(u(x),y)\)；我们有

\[
\Phi_ {u} (y, x) = \Phi (u (y), x) = \Phi (y, u (x)) = \overline {{\Phi (u (x) , y)}} = \overline {{\Phi_ {u} (x , y)}}.
\]

这表明 \(\Phi_{u}\) 是 E 上的埃尔米特型。反之，设 \(\Psi\) 是 E 上的一个埃尔米特型；由于从 E 到 E* 的映射 \(s_{\Phi}\)（它与 \(\Phi\) 相伴）是双射，故对每个 \(x \in \mathbf{E}\) ，存在 E 中唯一的元素 \(u(x)\) 使得 \(\Psi(x, y) = \Phi(u(x), y)\)；容易验证 \(u\) 是 E 的埃尔米特自同态。于是 \(u \to \Phi_{u}\) 是从 E 的埃尔米特自同态所成的集到 E 上埃尔米特型所成的集上的双射，称为典范双射。

设 A = K 或 A = K(i)；给定 E 上的一个双线性型 Ψ，对每个 \(x \in E\) ，存在唯一的元素 \(u(x)\)，即 E 中唯一的元素，使得 Ψ(x, y) = \(\overline{\Phi(u(x), y)}\)；立即可以看出 u 是 E 到其自身的（关于 J 的）半线性映射。由于 \(\overline{\Phi(u(x), y)} = \Phi(y, u(x)) = \overline{\Phi(u^{*}(y), x)}\)，可以看出，要使 Ψ 对称（相应地，交错），必须且只需 \(u^{*} = u\)（相应地，\(u^{*} = -u\)）。

定理 2. --- 设 S 是 E 的一族自同态（相应地，当 A = K 或 A = K(i) 时，为 E 到 E 中的半线性映射所成的集），它在映射 u → u* 下稳定。那么，若 V 是 E 的一个在 S 下稳定的子空间，则其正交补 V⁰ 也在 S 下稳定。此外，E 是在 S 下稳定、在 ≠ \{0\} 且在 S 下稳定的诸子空间中为极小的、且两两正交的一些子空间的直和。

实际上，设 V 是 E 的一个在 S 下稳定的子空间；对任意 \(x \in \mathrm{V}^0\) 、\(y \in \mathrm{V}\) 及 \(u \in \mathrm{S}\)，有 \(u^*(y) \in \mathrm{V}\)，从而 \(\Phi(y, u(x)) = \Phi(u^*(y), x) = 0\)（相应地，\(\Phi(y, u(x)) = \overline{\Phi(u^*(y), x)} = 0\)），于是 \(u(x) \in \mathrm{V}^0\)；因而 \(\mathrm{V}^0\) 在 S 下稳定，这就证明了第一个断言。对第二个断言，我们对 E 的维数 \(n\) 作归纳，\(n = 0\) 的情形是平凡的。当 \(n \neq 0\) 时，存在 E 的一个子空间 \(\mathrm{V} \neq \{0\}\)，它在 S 下稳定且极小，例如一个 \(\neq \{0\}\) 的、维数极小的稳定子空间。这时只需对 \(\mathrm{V}^0\) 应用归纳假设，因为 \(u\) 在 \(\mathrm{V}^0\) 上的限制（关于 \(\Phi\) 的限制）的伴随，等同于 \(\mathrm{V}^0\) 上 \(u\) 的伴随的限制。

推论 1. --- 设 A = K 或 A = K(i)。设 B 是 \(\mathfrak{L}_{\mathbf{A}}(\mathbf{E})\) 的一个在映射 \(u \to u^{*}\) 下稳定的子代数。则 E 是半单 B-模，且是两两正交的一些单子模的直和。代数 B 是半单的。

实际上，由于 E 的每个子 B-模 V 都有一个补，例如 \(V^{0}\)，故 B-模 V 是半单的（第八章 § 3 第 3 小节命题 7）。由于 E 的每个 \(\neq\{0\}\) 且极小的子 B-模都是单的，故 E 是两两正交的一些单子模的直和。最后，B 是半单代数，因为它有一个半单且忠实的模 E，其反向模是有限生成的（第八章 § 5 第 1 小节命题 3）。

推论 2. --- 在推论 1 的假设与记号下，再设代数 B 是交换的。则其所有元素都是 E 的（正规）半单自同态。当 A = K（相应地，A = K(i)）时，E 是两两正交的一些单 B-子模的直和，这些子模是 A 上维数为 1 或 2（相应地，1）的向量空间。

第一个断言由第八章 § 9 第 1 小节命题 2 得出，因为 B 是半单的。另一方面，每个单 B-模同构于形如 B/m 的 B-模，其中 m 是 B 的一个极大理想，从而同构于 A 上有限次的一个域 L；当 A = K（相应地，A = K(i)）时，域 L 必然同构于 K 或 K(i)（相应地，K(i)），因为 K(i) 是代数闭的（第六章 § 2 第 6 小节定理 3）。

命题 4. --- 设 u 是 E 的一个正规自同态。当 A 等于 K(i) 或 K 上的四元数体时，存在由 u 的特征向量组成的 E 的（关于 Φ 的）标准正交基。当 A = K 时，u 是半单的，且 E 是在 u 下稳定、两两正交、维数为 1 或 2 的一些子空间的直和。

先考察 A 交换的情形（A = K 或 A = K(i)）。这时子代数 B = A{[}u, u*{]}（ℒ\_A(E) 的）是交换的，因为 u 正规；由 § 1 第 8 小节的公式 (32) 与 (33)，它在映射 v → v* 下稳定。关于 A = K 情形的断言随即由定理 2 的推论 2 得出。当 A = K(i) 时，该推论还表明 E 是两两正交、在 u 下稳定、维数为 1 的一些向量子空间 Ax\_i (i = 1, . . ., n) 的直和；若令 \(e_{i} = (\Phi(x_{i}, x_{i}))^{-1/2}x_{i}\)，则 \((e_{i})\) 就是所要的标准正交基。

当 A 是 K 上的四元数体时，同样只需依定理 2 证明：在 u 与 u* 下稳定的、≠ \{0\} 的子空间所成的集的每个极小元素都具有维数 1。而今，这样的子空间 V 必含 u 的一个特征向量 x ≠ 0 (*)，这可如下看出：注意四元数体 A 含 K(i) 作为代数闭子域，并把 V 的标量域限制到 K(i)。于是只需证明 \(x\)（\(u\) 的特征向量）也是 \(u^{*}\) 的特征向量。设 \(u(x) = ax (a \in A)\)；这时我们有 \(\Phi(u(x), x) = a\Phi(x, x) = \Phi(x, x)a = \Phi(x, \bar{a}x)\)，因为 \(\Phi(x, x)\) 属于 A 的中心；而另一方面 \(\Phi(u(x), x) = \Phi(x, u^*(x))\)；由此得 \(\Phi(x, u^*(x) - \bar{a}x) = 0\)，因此可以写 \(u^*(x) = \bar{a}x + z\)，其中 \(z\) 是与 \(x\) 正交的一个向量。于是我们有

\[
\Phi (u ^ {*} (x), u ^ {*} (x)) = \bar {a} a \Phi (x, x) + \Phi (z, z) = \Phi (u (x), u (x)) + \Phi (z, z).
\]

而今，由于 \(u\) 是正规的，我们有

\[
\begin{array}{r l} \Phi (u (x), u (x)) & = \Phi (x, u ^ {*} u (x)) = \Phi (x, u u ^ {*} (x)) \\ & = \Phi (u u ^ {*} (x), x) = \Phi (u ^ {*} (x), u ^ {*} (x)). \end{array}
\]

于是有 \(\Phi(z, z) = 0\)，按假设这蕴含 \(z = 0\) 及 \(u^{*}(x) = \bar{a}x\)。证毕。

注记. --- 由定理 2 可知，对应于 \(u\) 的两个不同特征值的特征子空间是正交的。

命题 5. --- 设 u 是 E 的埃尔米特自同态。u 的特征值都属于 K，且存在由 u 的特征向量组成的 E 的标准正交基。

当 A 等于 K(i) 或 K 上的四元数体时，依命题 4，只需证明第一个断言；而今，若 x 是 u 的一个非零特征向量，a 是相应的特征值，则由假设 u = u*，我们有

\[
a \Phi (x, x) = \Phi (u (x), x) = \Phi (x, u (x)) = \Phi (x, x) \bar {a};
\]

由于 \(\Phi(x, x)\) 是 A 的中心中的一个非零元素，由此得 \(a = \bar{a}\)，从而 \(a \in K\)。因此，埃尔米特矩阵的所有特征值都在 K 中。现在设 A 等于 K。矩阵 \(M\)（即 \(u\) 关于 E 的一个标准正交基的矩阵）这时是对称的；于是，若把它看作 \(K(i)\) 上的矩阵，它就是一个埃尔米特矩阵。证明的第一部分于是表明，这个矩阵的所有特征值都在 K 中，从而当 \(E \neq \{0\}\) 时，它在 E 中有 \(\neq 0\) 的特征向量。由此推出，每个在 \(u\) 下稳定且极小的子空间必然具有维数 1，而结论如命题 4 中那样立即可得。

命题 6. --- a) 设 Ψ 是 E 上的一个埃尔米特型（相应地，当 A = K 或 A = K(i) 时，为对称双线性型）。存在 E 的一个基，它关于 Φ 是标准正交的，关于 Ψ 是正交的。

\begin{enumerate}
\def\labelenumi{\alph{enumi})}
\setcounter{enumi}{1}
\tightlist
\item
  设 A = K 或 A = K(i)，并设 Ψ 是 E 上的一个交错双线性型。存在 E 的一个基，它关于 Φ 是标准正交的，且 Ψ 关于这个基的矩阵具有以下形状
\end{enumerate}

\[
\left( \begin{array}{c c c c} 0 & a _ {1} & 0 & 0 \ldots 0 \\ - a _ {1} & 0 & 0 & 0 \ldots 0 \\ 0 & 0 & 0 & a _ {2} \ldots 0 \\ 0 & 0 & - a _ {2} & 0 \ldots 0 \\ \ldots \ldots \end{array} \right).
\]

其中诸 \(a_{i}\) 在 K 中 \(\geqslant 0\)。

当 \(\Psi\) 是埃尔米特型时，我们的断言由命题 5 以及 E 上埃尔米特型与关于 \(\Phi\) 的埃尔米特自同态之间的典范对应立即得出。在另外两种情形，设 \(u\) 是本小节开头由公式 \(\Psi(x, y) = \overline{\Phi(u(x), y)}\) 所定义的 E 到其自身的半线性映射。于是 \(u^2\) 是一个 A-线性映射；当 \(\Psi\) 对称（相应地，交错）时，有 \(u = u^*\)（相应地，\(u = -u^*\)），因而 \(u^2\) 是埃尔米特的；从而也有

\[
\begin{array}{r l} \Phi (u ^ {2} (x), x) & = \overline {{\Phi (x , u ^ {2} (x))}} = \Phi (u ^ {*} (x), u (x)) = \Phi (u (x), u (x)) \\ & (\text { resp. } - \Phi (u (x), u (x))) \end{array}
\]

这表明埃尔米特型 \((x, y) \to \Phi(u^{2}(x), y)\) 是正的（相应地，负的）。然后应用定理 2，设 V 是 E 的 \(\neq\{0\}\) 的、在 u 与 \(u^{*}\) 下稳定的子空间所成的族中的一个极小元素。由于 \(u^{2}\) 是埃尔米特自同态，V 含有一个 \(x \neq 0\)，它是 \(u^{2}\) 的特征向量；设 \(u^{2}(x) = ax\)，其中 \(a \in K\)（命题 5）。

当 \(\Psi\) 对称时，不等式 \(\Phi(u^{2}(x), x) \geqslant 0\) 表明 \(a \geqslant 0\)。设 \(y = a^{1/2}x + u(x)\)；我们有 \(u(y) = a^{1/2}u(x) + ax = a^{1/2}y\)，且 Ay 在 \(u\) 与 \(u^{*} = u\) 下稳定。若 \(y = 0\)，则 Ax 在 \(u\) 与 \(u^{*}\) 下稳定；在所有情形，V 都是维数为 1 的子空间，而我们的结论由定理 2 立即可得。

当 \(\Psi\) 交错时，有 \(a \leqslant 0\)；设

\[
y = (- a) ^ {1 / 2} x + u (x), \quad z = (- a) ^ {1 / 2} x - u (x).
\]

\[
\text { One has } u (y) = - (- a) ^ {1 / 2} z, u (z) = (- a) ^ {1 / 2} y \text { and }
\]

\[
\Phi (y, z) = - a \Phi (x, x) - \Phi (u (x), u (x)) = - \Phi (a x, x) + \Phi (u ^ {2} (x), x) = 0.
\]

若 \(y = z = 0\)，则有 \(a = 0\) 且 \(u(x) = 0\)，于是 \(Ax\) 在 \(u\) 与 \(u^{*}\) 下稳定，\(V\) 的维数为 1，且 \(\Psi\) 在 \(V\) 上的限制的矩阵为零，因为 \(\Psi\) 是交错的。否则，\(y\) 与 \(z\) 都 \(\neq 0\)，它们生成 \(V\)，且由于它们正交，\(V\) 的维数为 2，而 \(\Psi\) 在 \(V\) 上的限制的矩阵具有形状 \(\begin{pmatrix} 0 & b \\ -b & 0 \end{pmatrix}\)。最后，依公式 \(\Psi(x', y') = \overline{\Phi(u(x'), y')}\)，有 \(\Psi(x', y') = 0\)，只要 \(x'\) 与 \(y'\) 属于两个在 \(u\) 下稳定且关于 \(\Phi\) 正交的子空间。这就证明了存在 \(E\) 的（关于 \(\Phi\) 的）标准正交基，\(\Psi\) 关于它的矩阵具有所指出的形状。证毕。

习题. --- 1) 假设 § 7 开头的诸假设成立；设 \(\Phi\) 是 E 上的一个正埃尔米特型。

\begin{enumerate}
\def\labelenumi{\alph{enumi})}
\item
  要使 \(\Phi(x,x)\Phi(y,y)=\Phi(x,y)\overline{\Phi(x,y)}\)，必须且只需 \(x\) 与 \(y\) 线性相关，或由 \(x\) 与 \(y\) 生成的平面是迷向的。
\item
  假设对一切 \(x \in E\)，\(\Phi(x, x)\) 是 K 中的一个平方。证明对 E 中任意两个向量 x、y，我们有
\end{enumerate}

\[
\sqrt {\Phi (x + y , x + y)} \leqslant \sqrt {\Phi (x , x)} + \sqrt {\Phi (y , y)}.
\]

若 \(\Phi\) 非退化，则这个不等式的两边仅当 \(\alpha x + \beta y = 0\) 时才能相等，其中 \(\alpha\) 与 \(\beta\) 是 K 中两个不全为零的元素，且满足 \(\alpha \beta \leqslant 0\)。

\begin{enumerate}
\def\labelenumi{\arabic{enumi})}
\setcounter{enumi}{1}
\tightlist
\item
  假设 \(\mathbf{A} = \mathbf{K}\) 或 \(\mathbf{A} = \mathbf{K}(i)\)。设 \(X\) 是一个 \(n\) 阶埃尔米特方阵，元素在 \(\mathbf{A}\) 中，使得对每个非空子集 \(\mathbf{H}\) ⊂ \([1, n]\)，有 \(X_{\mathbf{H},\mathbf{H}} \geqslant 0\)（第 1 小节命题 3 的记号）。
\end{enumerate}

\begin{enumerate}
\def\labelenumi{\alph{enumi})}
\item
  设 \(\lambda\) 是 K 中一个 \(>0\) 的元素；证明矩阵 \(X + \lambda I\) 是正且非退化的（利用第 1 小节命题 3，对 \(n\) 作归纳论证）。
\item
  由此推出埃尔米特矩阵 X 是正的。
\end{enumerate}

\begin{enumerate}
\def\labelenumi{\arabic{enumi})}
\setcounter{enumi}{2}
\item
  假设 \(\mathbf{A} = \mathbf{K}\) 或 \(\mathbf{A} = \mathbf{K}(i)\)，且 E 是有限维的。设 \(\Phi_1, \Phi_2\) 是 E 上两个正埃尔米特型，\(V = (\alpha_{ij})\)、\(W = (\beta_{ij})\) 是这两个型关于 E 的同一基 \((e_i)\) 的矩阵。证明：埃尔米特型 \(\Phi\)（其矩阵 \((\gamma_{ij})\) 关于 \((e_i)\) 满足对每对指标 \(\gamma_{ij} = \alpha_{ij}\beta_{ij}\)）是正的；此外，若 \(\Phi_1\) 与 \(\Phi_2\) 非退化，则 \(\Phi\) 亦然。（在计算 \(\Phi(x, x)\) 时，把诸 \(\alpha_{ij}\) 用 \(\Phi_1(c_i, c_i)\) 的值表出，其中 \((c_i)\) 是 E 关于 \(\Phi_1\) 的一个正交基。）
\item
  设 A = K 或 A = K(i)。设 R 是 A 上一个 n 阶埃尔米特矩阵；称 R 具有符号差 \((s, t)\)，如果 \(A^{n}\) 上以 R 为其关于典范基的矩阵的埃尔米特型具有符号差 \((s, t)\)。以 \(\Delta_{k}\) 记从 R 中删去指标 \textgreater{} k 的行与列所得的 R 的主子式；假设 \(\Delta_{s+t} \neq 0\)，且对任何指标 \(k < s + t\)，\(\Delta_{k}\) 与 \(\Delta_{k+1}\) 不同时为零（参见 § 6 习题 1 b)）。证明：若 \(\Delta_{k} = 0\) 对某个 \(k < s + t\) 成立，则 \(\Delta_{k-1}\) 与 \(\Delta_{k+1}\) 符号相反（用 § 6 习题 1 a) 的方法），且数 s - t 等于
\end{enumerate}

\[
\operatorname{sgn} \Delta_ {1} + \operatorname{sgn} \left(\Delta_ {1} \Delta_ {2}\right) + \dots + \operatorname{sgn} \left(\Delta_ {s + t - 1} \Delta_ {s + t}\right)
\]

（利用 § 6 的习题 2）。

\begin{enumerate}
\def\labelenumi{\arabic{enumi})}
\setcounter{enumi}{4}
\item
  假设 A = K 或 A = K(i)；设 R、S 是 A 上两个埃尔米特矩阵，其符号差分别为 \((s, t)\) 与 \((s', t')\)（习题 4）；证明矩阵 \(R \otimes S\) 是符号差为 \((ss' + tt', st' + s't)\) 的埃尔米特矩阵。
\item
  设 K 是一个极大有序域，L 是 K 上一个有限维单代数。若 \((s, t)\) 是 L 上对称双线性型 \((x, y) \to \mathrm{Tr}_{\mathrm{L/K}}(xy)\) 的符号差，证明我们有：
\end{enumerate}

\(s-t=m\)，若 L 同构于 K 上一个 \(m\) 阶矩阵代数；

s - t = 0，若 L 同构于域 K(i) 上的矩阵代数；

\(s - t = -2m\)，若 L 同构于 K 上四元数体上的一个 \(m\) 阶矩阵代数。

\begin{enumerate}
\def\labelenumi{\arabic{enumi})}
\setcounter{enumi}{6}
\tightlist
\item
  假设 A 满足 § 7 开头的条件，且 E 具有有限维数 n。设 Φ 是 E 上的一个正且非退化的埃尔米特型。
\end{enumerate}

\begin{enumerate}
\def\labelenumi{\alph{enumi})}
\item
  证明每个关于 \(\Phi\) 的相似变换都可以以唯一的方式写成一个 K 中比 \(\geqslant 0\) 的位似与一个酉变换之积。
\item
  对 E 的每个基 \((a_{i})_{1\leqslant i\leqslant n}\)，证明存在一个且只有一个 E 的标准正交基 \((e_{i})_{1\leqslant i\leqslant n}\) 满足下列条件：\(1^{\circ}\) 对每个 \(m\)（\(1\leqslant m\leqslant n\)），由 \(a_1,\ldots ,a_m\) 生成的子空间等同于由 \(e_1,\ldots ,e_m\) 生成的子空间；\(2^{\circ}\) 在 K 中有 \(\Phi (a_i,e_i) > 0\)（对每个指标 \(i\)）（参见 § 6 第 1 小节命题 1）。
\item
  由 \(b\) 推出：对 A 上每个可逆方阵 \(M\)（\(n\) 阶），存在一个且只有一个矩阵对 \((L, U)\)（\(n\) 阶），使得 \(U\) 是酉矩阵，\(L = (\lambda_{ij})\) 的对角线下方全为零，且对角元素 \(\lambda_{ii}\) 属于 K 并 \(>0\)，最后 \(M = LU\)。
\end{enumerate}

设 A = K；设 L 是 Λ 上一个有限维埃尔米特空间，其度量型是正且非退化的。设 M、N 是 L 的两个子集，u 是从 M 到 N 上的一个双射，使得对 M 中任意点 a、b 有 \(e(u(a), u(b)) = e(a, b)\)，证明存在一个位移，它在 M 上的限制等于 u（按格拉姆–施密特正交化方法的方式，对由 M 生成的仿射线性簇的维数作归纳论证）。这个命题能否推广到 A 等于 K(i) 或 K 上四元数体的情形？

\begin{enumerate}
\def\labelenumi{\arabic{enumi})}
\setcounter{enumi}{8}
\item
  假设第 3 小节的诸条件成立，此外 A 是 K 上的四元数体。设 u 是 E 的一个正规自同态。证明 E 是两两正交的子空间 \(F_{k}\)（\(1 \leqslant k \leqslant r\)）的直和，使得在每个 \(F_{k}\) 中存在一个标准正交基（\(e_{ik}\)）（\(1 \leqslant i \leqslant n_{k}\)），且 \(u(e_{ik}) = \lambda_{k}e_{ik}\)（\(1 \leqslant i \leqslant n_{k}\)），而指标不同的两个 \(\lambda_{k}\) 不能通过 A 的一个内自同构互相变换。此外，若（\(F_{k}'\)）是 E 的一个具有同样性质的第二个分解，带有特征值组（\(\lambda_{k}'\)），则有 \(F_{k}' = F_{k}\)（在指标的一个置换下），且 \(\lambda_{k}' = \alpha_{k}\lambda_{k}\alpha_{k}^{-1}\)；u 关于特征值 \(\lambda_{k}\) 的特征向量所成的集合是 \(A_{k}\)（\(\lambda_{k}\) 在 A 中的中心化子）上由诸 \(e_{ik}\)（\(1 \leqslant i \leqslant n_{k}\)）生成的子空间。
\item
  假设第 3 小节的诸条件成立，此外 A 等于 K(i) 或 K 上的四元数体。设 u 是 E 的一个正规自同态。
\end{enumerate}

\begin{enumerate}
\def\labelenumi{\alph{enumi})}
\item
  证明：要使 E 的一个向量子空间 F 满足 \(u(\mathrm{F}) \subset \mathrm{F}\)，必须且只需 F 由 u 的特征向量生成；此时我们有 \(u(\mathrm{F}^{0}) \subset \mathrm{F}^{0}\) 和 \(u^{*}(\mathrm{F}) \subset \mathrm{F}\)。
\item
  要使 \(u\) 是酉的（相应地，\(u^{*} = u\)、\(u^{*} = -u\)），必须且只需对每个特征值 \(\lambda\)（\(u\) 的特征值）有 \(\lambda \bar{\lambda} = 1\)（相应地，\(\bar{\lambda} = \lambda\)、\(\bar{\lambda} = -\lambda\)）。
\item
  我们称 E 的埃尔米特自同态 \(u\) 是正的（相应地，正且非退化的），如果典范地对应于它的埃尔米特型 \((x, y) \to \Phi(u(x), y)\) 是正的（相应地，正且非退化的）；为此必须且只需 \(u\) 的所有特征值都 \(\geqslant 0\)（相应地，\(>0\)）。
\item
  证明对每个整数 m\textgreater{} 0，存在 E 的一个正规自同态 v 使得 \(v^{m}=u\)。若 u 是正埃尔米特的，则存在唯一的正埃尔米特自同态 v 使得 \(v^{m}=u\)，且存在一个多项式 \(f\in K[X]\) 使得 \(v=f(u)\)；这最后一个结果当 A = K 时也成立。
\end{enumerate}

\begin{enumerate}
\def\labelenumi{\arabic{enumi})}
\setcounter{enumi}{10}
\tightlist
\item
  假设第 3 小节的诸条件成立，此外 A = K。设 u 是 E 的一个正规自同态。
\end{enumerate}

\begin{enumerate}
\def\labelenumi{\alph{enumi})}
\item
  设 V 是 E 的不退缩为 0 且在 u 与 \(u^{*}\) 下稳定的子空间所成之集的一个极小元素；证明若 V 的维数为 2，则 u 在 V 上的限制是一个乘子 \textgreater{} 0 的直接相似。
\item
  证明 E 的每个在 u 下稳定的子空间也在 \(u^{*}\) 下稳定。（设 \(E_{0}\) 是把 E 的纯量域扩张到 \(\mathrm{K}(i)\) 后得到的向量空间；\(\Phi\) 是 \(E_{0}\) 上一个正非退化埃尔米特型在 E 上的限制，而 u 是 \(u_{0}\)（\(E_{0}\) 的一个对这个型而言的正规自同态）在 E 上的限制；此外，E 是满足下列条件的 \(x \in E_{0}\) 所成的集合：x 在 \(E_{0}\) 的一个对合半线性双射 j 下不变，且我们有 \(u_{0}j = ju_{0}\)。然后应用习题 10 a)。）
\end{enumerate}

\begin{enumerate}
\def\labelenumi{\arabic{enumi})}
\setcounter{enumi}{11}
\item
  假设第 3 小节的诸条件成立，此外 A 等于 K(i) 或 K 上的四元数体。证明：对 E 的每个自同态 u，存在 E 的一个标准正交基，使 u 关于这个基的矩阵在对角线下方只有零。（对 E 的维数作归纳论证，考虑 u 的一个特征向量。）
\item
  设 A 是一个满足 § 7 开头条件的体，E、F 是 A 上两个有限维向量空间，\(\Phi\)（相应地，\(\Psi\)）是 E（相应地，F）上的一个正非退化埃尔米特型。证明：对每个从 E 到 F 的线性映射 \(u\)，其伴随 \(u^{*}\)（\(u\) 的伴随，对 \(\Phi\) 与 \(\Psi\) 而言；参见 § 1 第 8 小节）具有如下性质：\(u^{*}u\) 与 \(uu^{*}\) 分别是 E 与 F 的正埃尔米特自同态（习题 10 c)）。
\item
  假设第 3 小节的诸条件成立。设 u 是 E 的一个自同态，\(h_{1}\)、\(h_{2}\) 是 E 的两个正埃尔米特自同态，使得我们有 \(h_{1}^{2}=u^{*}u\)、\(h_{2}^{2}=uu^{*}\)（习题 13 与 10 d)）。
\end{enumerate}

\begin{enumerate}
\def\labelenumi{\alph{enumi})}
\item
  证明存在一个酉自同态 \(\nu\) 使得 \(u = \nu h_{1} = h_{2}\nu\)，特别地 \(h_{1}\) 与 \(h_{2}\) 相似（注意 \(\bar{u}(0) = \bar{h}_{1}^{-1}(0)\)，且若 V 是与 \(\bar{u}(0)\) 正交的子空间，则 \(\Phi(u(x), u(x)) = \Phi(h_{1}(x), h_{1}(x))\) 对一切 \(x \in \mathrm{V}\) 成立。要使 \(\nu\) 被唯一确定，必须且只需 \(u\) 是双射的。要能取 \(\nu\) 与 \(h_{1}\) 交换，必须且只需 \(u\) 是正规的。
\item
  由 a) 推出：A 上每个 n 阶方阵 M 都可以写成 UDV，其中 U 与 V 是酉矩阵，D 是一个对角矩阵，其元素 \(\geqslant 0\) 且以 MM* 的特征值作为其平方。
\end{enumerate}

\begin{enumerate}
\def\labelenumi{\arabic{enumi})}
\setcounter{enumi}{14}
\tightlist
\item
  假设第 3 小节的诸条件成立。证明 A 上每个正埃尔米特矩阵 H 都可以写成 \(LL^{*}\) 的形式，其中 \(L = (\lambda_{ij})\) 的对角线下方只有零，且对角线项属于 K 并 \(\geqslant 0\)；此外，当 H 可逆时 L 由这些条件唯一确定（参见习题 10 d) 与 7 c)）。
\end{enumerate}

¶ 16) 假设第 3 小节的诸条件成立，此外 A = K(i)。设 u 是 E 的一个自同态。

\begin{enumerate}
\def\labelenumi{\alph{enumi})}
\item
  \(u\) 的特征值所成的集合含于 \(\Phi(x, u(x))\) 当 \(x\) 遍历 E 中满足 \(\Phi(x, x) = 1\) 的元素所成之集时所取值的集合 U。
\item
  我们称 A = K(i) 的一个子集 C 是凸的，如果对每对元素 \((\xi, \eta) \in \mathrm{C}^{2}\) 和每个 \(\tau \in K\)（满足 \(0 \leqslant \tau \leqslant 1\)），有 \(\tau\xi + (1 - \tau)\eta \in C\)。证明集合 U 是凸的。（化归到 n = 2 的情形；把 u 写成 \(\nu + iw\) 的形式，其中 \(\nu\) 与 w 是埃尔米特的，并在必要时把 u 换成 \(\lambda u\)（其中 \(\lambda \in A\)，\(\lambda\overline{\lambda} = 1\)），证明一切都化归为证明下列性质。设 \(f(\xi_{1}, \xi_{2}) = \xi_{1}\overline{\xi}_{1} + \xi_{2}\overline{\xi}_{2}\)、\(g(\xi_{1}, \xi_{2}) = a\xi_{1}\overline{\xi}_{1} + b\xi_{2}\overline{\xi}_{2}\)（\(a \in K, b \in K\)）、\(h(\xi_{1}, \xi_{2}) = \alpha\xi_{1}\overline{\xi}_{1} + \beta\xi_{1}\overline{\xi}_{2} + \beta\xi_{2}\overline{\xi}_{1} + \gamma\xi_{2}\overline{\xi}_{2}\)（\(\alpha \in K, \gamma \in K, \beta \in A\)）；设 \((\eta_{1}, \eta_{2}) \in A^{2}\)、\((\zeta_{1}, \zeta_{2}) \in A^{2}\) 满足 \(f(\eta_{1}, \eta_{2}) = f(\zeta_{1}, \zeta_{2}) = 1\)、\(g(\eta_{1}, \eta_{2}) = g(\zeta_{1}, \zeta_{2}) = 1\)、\(h(\eta_{1}, \eta_{2}) > 0\)、\(h(\zeta_{1}, \zeta_{2}) < 0\)；则存在 \((\theta_{1}, \theta_{2}) \in A^{2}\) 满足 \(f(\theta_{1}, \theta_{2}) = 1\)、\(g(\theta_{1}, \theta_{2}) = 1\) 与 \(h(\theta_{1}, \theta_{2}) = 0\)。可以先指出，对每对 \((\xi_{1}, \xi_{2})\)，存在 \(\mu \in A\) 满足 \(\mu\overline{\mu} = 1\) 与 \(\beta\mu\xi_{1}\overline{\xi}_{2} + \beta\overline{\mu}\xi_{2}\overline{\xi}_{1} = 0\)，这将使我们可以化归到 \(\beta = 0\) 的情形；利用第六章 § 2 第 5 小节命题 5。）
\item
  证明若 \(u\) 是正规的，则 U 是包含 \(u\) 的所有特征值的最小凸集。给出一个 \(u\) 不是正规的但 U 仍具有上述性质的例子。（取 \(u\) 的特征值为元素 \(\pm 1 \pm i\)，它们是单特征值，并取 0 作为一个二重特征值。）
\end{enumerate}

¶ 17) 我们假设第 3 小节的诸条件成立；对每个 \(\xi\in\mathbf{A}\)，我们用 \(|\xi|\) 记 K 的元素 \(\rho\geqslant0\)，它满足 \(\rho^{2}=\xi\bar{\xi}\)（\(\xi\) 的绝对值）。对 A 上每个方阵 \(M=(\alpha_{ij})\)，我们置 \(f(M)=\max_{i}|\alpha_{ii}|\)、\(g(M)=\max_{i,j}|\alpha_{ij}|\)，并用 \(\varphi(M)\) 记 \(M\) 的特征值的最大绝对值（将证明当 A 是 K 上的四元数体时这个定义是有意义的）。

\begin{enumerate}
\def\labelenumi{\alph{enumi})}
\tightlist
\item
  设 \(A, B, D\) 是 A 上三个 \(n\) 阶方阵。证明：若 \(D\) 是对角矩阵，则
\end{enumerate}

\[
g ^ {2} (A D B ^ {*}) \leqslant f ^ {2} (D) f (A A ^ {*}) f (B B ^ {*})\tag{1}
\]

（利用第 1 小节的不等式 (1)）。由此推出

\[
g (A B B ^ {*} A ^ {*}) \leqslant f (A A ^ {*}) \varphi (B B ^ {*})\tag{2}
\]

（把命题 5 应用于埃尔米特矩阵 \(BB^{*}\)）。由此推出：对 A 上 \(m\) 个任意方阵 \(A_{i}\)（\(1 \leqslant i \leqslant m\)），我们有

\begin{enumerate}
\def\labelenumi{(\arabic{enumi})}
\setcounter{enumi}{2}
\tightlist
\item
\end{enumerate}

\[
g ^ {2} (A _ {1} A _ {2} \dots A _ {m}) \leqslant f (A _ {1} A _ {1} ^ {*}) \varphi (A _ {2} A _ {2} ^ {*}) \dots \varphi (A _ {m - 1} A _ {m - 1} ^ {*}) f (A _ {m} ^ {*} A _ {m})\tag{4}
\]

\[
\varphi^ {2} (A _ {1} A _ {2} \dots A _ {m}) \leqslant \varphi (A _ {1} A _ {1} ^ {*}) \varphi (A _ {2} A _ {2} ^ {*}) \dots \varphi (A _ {m} A _ {m} ^ {*})\tag{5}
\]

\[
g ^ {2} (A _ {1} A _ {2} \dots A _ {m}) \leqslant \varphi (A _ {1} A _ {1} ^ {*}) \dots \varphi (A _ {m} A _ {m} ^ {*}).
\]

（对 (3)，从 (2) 出发用归纳法论证。对 (4)，利用习题 12；最后由 (3) 推出 (5)。）

  证明当把 \(A_{m}^{*}A_{m}\) 换成 \(A_{m}A_{m}^{*}\) 时（注意可以有 \(f(A^{*}A)\neq f(A A^{*})\)），或者当把 \(\varphi(A_{i}A_{i}^{*})\) 换成 \(f(A_{i}A_{i}^{*})\)（对 \(2\leqslant i\leqslant m-1\)）时（取所有 \(A_{i}\) 都等于全元素为 1 的方阵），不等式 (3) 不再必然成立。

\begin{enumerate}
\def\labelenumi{\alph{enumi})}
\setcounter{enumi}{1}
\tightlist
\item
  设 u 是 E 的一个正规自同态；若 λ 是 u 的一个特征值（当 A = K 时，是 K(i) 的一个元素），则 \(\lambda\bar{\lambda}\) 是 u*u 的一个特征值。因此，对每个正规矩阵 M（E 的一个正规自同态关于一个标准正交基的矩阵），我们有 φ²(M) = φ(MM*)。由此推出也有 g(M) ≤ φ(M)，且更一般地，若 M₁, . . ., Mₘ 是正规的，则
\end{enumerate}

\begin{enumerate}
\def\labelenumi{(\arabic{enumi})}
\setcounter{enumi}{5}
\tightlist
\item
\end{enumerate}

\[
g (M _ {1} \dots M _ {m}) \leqslant \varphi (M _ {1}) \dots \varphi (M _ {m})\tag{7}
\]

\[
\varphi (M _ {1} \dots M _ {m}) \leqslant \varphi (M _ {1}) \dots \varphi (M _ {m})
\]

（利用 (4) 与 (5)）。

\begin{enumerate}
\def\labelenumi{\alph{enumi})}
\setcounter{enumi}{2}
\tightlist
\item
  证明若 \(H\) 是 A 上一个正埃尔米特矩阵，我们有 \(f(H) = g(H) \leqslant \varphi(H)\)（利用 (3) 与 b)，把 H 写成 \(H = AA^{*}\)）。
\end{enumerate}

¶ 18) 我们假设第 3 小节的诸条件成立。

\begin{enumerate}
\def\labelenumi{\alph{enumi})}
\tightlist
\item
  对 E 的每个自同态 u，设 \((e_{i})\) 是由 \(u^{*}u\) 的特征向量组成的一个 E 的标准正交基，\(\rho_{i}\) 是 \(u^{*}u\) 对应于向量 \(e_{i}\) 的特征值；我们置 \(s(u) = (\sum_{i=1}^{n} \rho_{i})^{1/2}\)（当 A 交换时，这是 \(\mathrm{Tr}(u^{*}u)\) 的平方根）；我们有 \(s(u^{*}) = s(u)\)。若 u、v 是 E 的两个自同态，U、V 是它们关于 E 的同一标准正交基的矩阵，证明对元素在 K 中的矩阵 \(UV^{*} + VU^{*}\)，有 \((\mathrm{Tr}(UV^{*} + VU^{*}))^{2} \leqslant 4s(u)s(v)\)（注意 \((u^{*} + \lambda v^{*})(u + \lambda v)\) 对每个 \(\lambda \in K\) 是正埃尔米特的）；由此推出 \(s(u + v) \leqslant s(u) + s(v)\)。若此外 A 是交换的，则有
\end{enumerate}

\[
\left| \operatorname{Tr} (u v) \right| ^ {2} \leqslant s (u) s (v) \quad \text { and } \quad \left| \operatorname{Tr} (u) \right| \leqslant \sqrt {n} s (u)
\]

（利用习题 14 b) 把 \(u\) 写成一个乘积）。

\begin{enumerate}
\def\labelenumi{\alph{enumi})}
\setcounter{enumi}{1}
\tightlist
\item
  我们进一步假设 A 是交换的（故等于 K 或 K(i)）。若 \(H_{1}, H_{2}\) 是两个正埃尔米特方阵，证明有 \(\mathrm{Tr}(H_{1}H_{2}) \geqslant 0\)（化归到 \(H_{2}\) 是对角矩阵的情形）。由此推出我们有（用习题 17 的记号）
\end{enumerate}

\[
\left| \operatorname{Tr} (H _ {1} H _ {2}) \right| \leqslant \varphi (H _ {1}) \operatorname{Tr} (H _ {2}) \leqslant \operatorname{Tr} (H _ {1}) \operatorname{Tr} (H _ {2}).
\]

  从这些不等式推出：对 E 的任意两个自同态 \(u, \nu\)，我们有 \(s(u\nu) \leqslant s(u)s(\nu)\)。

\begin{enumerate}
\def\labelenumi{\alph{enumi})}
\setcounter{enumi}{2}
\tightlist
\item
  仍假设 A 是交换的，设 \(\prod_{i=1}^{n}(x-\lambda_i)\) 是 E 的任一自同态 \(u\) 的特征多项式的分解为线性因式之积；证明有 \(\sum_{i=1}^{n}|\lambda_i|^2\leqslant(s(u))^2\)，且要使不等式两边相等，必须且只需 \(u\) 是正规的（利用习题 12）。
\end{enumerate}

\begin{enumerate}
\def\labelenumi{\arabic{enumi})}
\setcounter{enumi}{18}
\item
  假设第 3 小节的诸条件成立。设 M 是 A 上一个 n 阶方阵；证明对 M 的每个方子矩阵 N（从 M 中删去若干行以及与这些行同指标的列而得到），有（用习题 17 的记号）\(\varphi^{2}(N) \leqslant \varphi(MM^{*})\)（适当地应用习题 17 的公式 (4)）。特别地，若 M 是正规矩阵，则 \(\varphi(N) \leqslant \varphi(M)\)（参见习题 17 b)）。
\item
  我们假设第 3 小节的诸条件成立，此外 A 等于 K 或 K(i)。把习题 17 至 19 的结果应用于 Φ 到 p 次外幂的扩张（§ 1 第 9 小节）以及所考虑的自同态或矩阵的 p 次外幂。特别地，证明：若 \(\prod_{i=1}^{n}(X-\lambda_i)\) 与 \(\prod_{i=1}^{n}(X-\rho_i^2)\) 分别是某个
\end{enumerate}

E 的自同态 \(u\) 与自同态 \(u^*u\) 的特征多项式的线性因式分解，且若假设 \(|\lambda_i| \geqslant |\lambda_{i+1}|\) 与 \(\rho_i \geqslant \rho_{i+1} \geqslant 0\)（\(1 \leqslant i \leqslant n-1\)），则有

\[
\mid \lambda_ {1} \lambda_ {2} \dots \lambda_ {h} \mid \leqslant \rho_ {1} \rho_ {2} \dots \rho_ {h}
\]

对 \(1 \leqslant h \leqslant n - 1\) 成立，且有 \(|\lambda_1\lambda_2 \ldots \lambda_n| = \rho_1\rho_2 \ldots \rho_n\)。

\begin{enumerate}
\def\labelenumi{\arabic{enumi})}
\setcounter{enumi}{20}
\tightlist
\item
  我们假设第 3 小节的诸条件成立。设 u 是 E 的一个埃尔米特自同态，\(\prod_{i=1}^{n}(X-\lambda_{i})\) 是它的特征多项式的线性因式分解；我们假设 \(\lambda_{i}\geqslant\lambda_{i+1}\)（\(1\leqslant i\leqslant n-1\)）。
\end{enumerate}

\begin{enumerate}
\def\labelenumi{\alph{enumi})}
\item
  证明 K 中的特征值 \(\lambda_{i}\) 中最大者（相应地，最小者）等于 \(\Phi(u(x), x)\) 的值中最大者（相应地，最小者），这里 \(x\) 遍历由这样的 \(x \in E\) 所成的集合：\(\Phi(x, x) = 1\)。（直接论证，或者化归到 A = K(i) 的情形后应用习题 16 c)（ \(\S\) 3 习题 4）。）
\item
  设 \(\Psi_{\mathrm{v}}\) 是埃尔米特型 \(\Psi\)（与 \(u\) 相伴的）在 E 的一个向量子空间 V 上的限制，\(u_{\mathrm{v}}\) 是与 \(\Psi_{\mathrm{v}}\) 相伴的 V 的埃尔米特自同态。证明 \(\lambda_{k}\) 是诸 \(u_{\mathrm{v}}\)（当 V 遍历 E 的维数为 \(n - k + 1\) 的向量子空间所成之集时）的最大特征值中最小的一个（利用命题 5）。
\end{enumerate}

\begin{enumerate}
\def\labelenumi{\arabic{enumi})}
\setcounter{enumi}{21}
\tightlist
\item
  我们假设第 3 小节的诸条件成立，此外 A 等于 K(i) 或 K 上的四元数体。设 u、v 是 E 的两个正规自同态；设 (E \(_{i}\) ) \(_{1\leqslant i\leqslant r}\)（相应地，(F \(_{j}\) ) \(_{1\leqslant j\leqslant s}\)）是 E 分解为两两正交的子空间之直和，使得在 E \(_{i}\)（相应地，F \(_{j}\)）中存在一个由 u（相应地，v）的属于同一特征值 \(\lambda_{i}\)（相应地，\(\mu_{j}\)）的特征向量组成的标准正交基，且 \(\lambda_{h}\) 与 \(\lambda_{i}\)（相应地，\(\mu_{j}\) 与 \(\mu_{k}\)）当 \(h\neq i\)（相应地，\(j\neq k\)）时不能通过 A 的一个内自同构互相变换（参见命题 4 与习题 9）。要使 E 的一个自同态 w 满足 uw = wν，必须且只需对每个 j（\(1\leqslant j\leqslant s\)），w(F \(_{j}\) ) 含于某个 E \(_{i}\) 之中，且 w 使 v 的属于特征值 \(\mu_{j}\) 的每个特征向量的像是 u 的属于特征值 \(\mu_{j}\) 的一个特征向量（特别地，这蕴涵 \(\mu_{j}\) 与 \(\lambda_{i}\) 可通过 A 的一个内自同构互相变换）。由此推出：若确实如此，则有 u*w = wν*。
\end{enumerate}

¶ 23) 我们假设第 3 小节的诸条件成立。设 u 与 v 是 E 的两个自同态。

\begin{enumerate}
\def\labelenumi{\alph{enumi})}
\item
  我们假设 u 与 uv 是正规的。要使 vu 是正规的，必须且只需 v 与 \(u^{*}u\) 交换。（要看出条件是必要的，利用关系式 \(u(vu) = (uv)u\) 与习题 22；要看出条件是充分的，利用习题 14 a) 与 10 d)。）
\item
  我们假设 \(u, v\) 与 \(uv\) 正规；设 \(\beta\) 是 \(uu^*\) 的最大特征值，\(F\) 是 \(E\) 的由 \(uu^*\) 的属于特征值 \(\beta\) 的特征向量所成的子空间。证明 \(v(F) \subset F\)。（注意对每个正规自同态 \(u\) 和每个 \(x \in E\)，有
\end{enumerate}

\[
\Phi (u (x), u (x)) = \Phi (u ^ {*} (x), u ^ {*} (x)).
\]

  化归到 A 等于 K(i) 或 K 上的四元数体的情形；于是 F 具有一个由 u（与 \(u^{*}\)）的特征向量组成的基；注意对这样的向量 z，我们有 \(\Phi(u^{*}u\nu(z),\nu(z))=\beta\Phi(\nu(z),\nu(z))\)。由此推出 \(\nu u\) 是正规的（对 \(u^{*}u\) 的不同特征值的个数作归纳论证）。若 h（相应地，\(h'\)）是使 \(h^{2}=uu^{*}\)（相应地，\(h'^{2}=\nu\nu^{*}\)）的正埃尔米特自同态，且若置 \(u=hu_{1}\)、\(\nu=h'\nu_{1}\)，其中 \(u_{1}\) 与 \(\nu_{1}\) 是酉的，h 与 \(u_{1}\) 交换且 \(h'\) 与 \(\nu_{1}\) 交换（习题 14 a)），证明偶 \((h,h')\)、\((h,\nu_{1})\) 与 \((h',u_{1})\) 都是可交换的；反之亦然。由此推出 \(u^{m}\nu^{n}\)、\(\nu^{n}u^{m}\)、\(u\nu^{*}\) 与 \(\nu^{*}u\) 这时都是正规的（m 与 n 是任意整数 \textgreater0）。

¶ 24) 假设第 3 小节的诸条件成立。设 Γ 是 E 的一个自同构群，使得每个 \(u \in \Gamma\) 都是正规的。证明存在 E 的一个分解，分解为两两正交的子空间 \(E_k\)（\(1 \leqslant k \leqslant r\)）的直和，使得 \(E_k\) 上每个 \(u \in \Gamma\) 的限制都具有 \(\lambda_k v_k\) 的形式，其中 \(\lambda_k\) 是 K 的一个元素 \textgreater0，\(v_k\) 是 \(E_k\) 的一个酉自同态（把每个 \(u \in \Gamma\) 分解成 \(h\nu\) 的形式，其中 \(v\) 是酉的，\(h\) 是正埃尔米特的，且 \(h^2 = uu^*\)（习题 14a)；利用习题 23b 并把定理 2 的推论 2 应用于由埃尔米特自同态 \(h\)（它们对应于诸 \(u \in \Gamma\)）生成的代数）。由此推出：若 Γ 是有限的，则诸元素 \(u \in \Gamma\) 都是酉的。

\begin{enumerate}
\def\labelenumi{\arabic{enumi})}
\setcounter{enumi}{24}
\tightlist
\item
  假设 A = K（极大有序域）。设 L 是 A 上一个 n 维欧几里得空间，其度量型是正且非退化的。
\end{enumerate}

  设 S 是 L 中一个非退化仿射二次曲面（§ 6 习题 25）。若 S 具有一个中心 a，且在 L 中取 a 为原点，证明存在 L 的一个标准正交基 \((e_{i})\)，使得关于这个基，S 有方程 \(\lambda_{1}\xi_{1}^{2} + \cdots + \lambda_{n}\xi_{n}^{2} = 1\)。此外，若两个标准正交基都具有这个性质，则与它们对应的元素 \(\lambda_{i} \in K\) 除次序外是相同的。

  若 S 没有中心，证明存在 S 的一个点 b 以及（取 b 为原点）L 的一个标准正交基，使得关于这个基，S 有方程 \(\lambda_{1}\xi_{1}^{2} + \cdots + \lambda_{n-1}\xi_{n-1}^{2} + \xi_{n} = 0\)。（若 \(\overline{S}\) 是使 \(S = L \cap \overline{S}\) 的射影二次曲面，c 是关于 \(\overline{S}\) 的无穷远超平面 \(H_{0}\) 的（无穷远）极点，则由下列条件确定 b：过 b 与 c 的直线垂直于 S 在点 b 处的切超平面。）

¶ 26) a) 设 K 是一个交换域，S 是 K 的一个子集，使得 K 是极大 S 可序域（第六章 § 2 习题 8）。设 f 是 K{[}X{]} 中一个多项式，L 是它的分裂域。证明：若对 K 上每个与其环结构相容的（全）序结构，L 都容许 K 的一个有序扩张结构，则必有 L = K。（用反证法论证：按第六章 § 2 习题 8 e) 的方式进行，化归到将有 {[}L : K{]} = 2 的情形。最后注意，若 b ∈ K 不是平方，则存在 K 的一个与其环结构相容的全序结构，使 b \textless{} 0（参见第六章 § 2 第 3 小节定理 1 的引理）。）

\begin{enumerate}
\def\labelenumi{\alph{enumi})}
\setcounter{enumi}{1}
\tightlist
\item
  把第 3 小节的结果（命题 6 除外）推广到 K 是极大 S 可序域的情形 (*)。（先利用 a) 证明命题 5 的类似命题。然后对交换代数 B 由埃尔米特自同态组成的情形，利用第八章 § 9 第 4 小节命题 10，建立定理 2 推论 2 的类似命题。对 A = K(i) 时正规自同态 u 的情形，注意这时可以把 u 写成 u = v + iw，其中 v 与 w 是埃尔米特的且交换。最后，若 A 是 K 上的四元数体，E₀ 是把 E 看作 K(i) 上维数 2n 的向量空间所得的空间，且若置 Φ(x, y) = Φ₁(x, y) + Φ₂(x, y)j，其中 Φ₁ 与 Φ₂ 取值在 K(i) 中，则注意若 u 对 Φ 是正规的，它对 Φ₁ 也是正规的。）
\end{enumerate}

¶ 27) 设 K 是一个极大有序域，E 是 K 上一个 n 维向量空间，Φ 是 E 上一个非退化对称双线性型，其符号差 (s, t) 不同于 (n, 0) 与 (0, n)。设 (ei) 是 Φ 的一个正交基，使得对 1 ≤ i ≤ s 有 Φ(ei, ei) = 1，对 s + 1 ≤ i ≤ n 有 Φ(ei, ei) = -1。对每个正交变换 u ∈ U(Φ)，设

\[
U = \left( \begin{array}{c c} M & N \\ P & Q \end{array} \right)
\]

\(u\) 关于 \((e_{i})\) 的矩阵，把它写成矩阵的方块阵列的形式，对应于把 \([1, n]\) 分为 \([1, s]\) 与 \([s + 1, n]\) 的分划。

\begin{enumerate}
\def\labelenumi{\alph{enumi})}
\tightlist
\item
  证明下列关系式
\end{enumerate}

\[
\begin{array}{r l} ^ {t} M. M - ^ {t} P. P & = I _ {s} \\ ^ {t} Q. Q - ^ {t} N. N & = I _ {t} \\ ^ {t} M. N - ^ {t} P. Q & = 0. \end{array}
\]

\begin{enumerate}
\def\labelenumi{\alph{enumi})}
\setcounter{enumi}{1}
\item
  设 R 是 K 上一个 t 行 s 列矩阵，使得 \(I_{s} - ^{t}R \cdot R\) 是一个正非退化埃尔米特型（在 \(K^{s}\) 上）的矩阵。证明 det \((M + \lambda NR)\) 在 K 中对 \(-1 \leqslant \lambda \leqslant 1\) 不变号（利用 a)，证明 \(^{t}(M + \lambda NR)(M + \lambda NR)\) 是一个正非退化对称型的矩阵）。
\item
  置 \(\sigma(u) = \sigma(U) = \operatorname{sgn}(\det M)\)；证明对任意两个元素 \(u\)、\(\nu\)（属于正交群 \(\mathbf{U}(\Phi)\)），有 \(\sigma(u\nu) = \sigma(u)\sigma(\nu)\)（利用 \(a\) 与 \(b\)）。由此推出特殊正交群 \(\mathbf{SU}(\Phi)\) 的换位子群不同于 \(\mathbf{SU}(\Phi)\)（参见 § 10 习题 9）。
\end{enumerate}

\section{二次型的型}

在本节中我们假设 A 是一个交换域。

\subsection{二次型的型.}

  给定 A 上向量空间 E 上的一个二次型 Q（§ 3 第 4 小节），我们称 E 为 Q 的定义空间，称 dim(E) 为 Q 的维数。给定 A 上向量空间 E、E' 上的两个二次型 Q、Q'，我们用 Q τ Q' 记它们的直和（§ 3 第 4 小节）。回忆，两个中性型的直和是中性的（§ 4 第 2 小节）。

  引入下列关系：

「Q 与 Q' 是 A 上有限维的非退化二次型，且存在中性二次型 N、N'，使得 Q τ N 等价于 Q' τ N'」。

  这个关系，我们将用 \(Q \sim Q'\) 来记，显然是自反的且对称的。它还是传递的：事实上，若 \(Q\)、\(Q'\)、\(Q''\) 是使 \(Q \sim Q'\) 与 \(Q' \sim Q''\) 成立的二次型，则存在中性二次型 \(M\)、\(M'\)、\(N\)、\(N'\)，使得 \(Q \tau M\) 等价于 \(Q' \tau M'\)，且 \(Q' \tau N\) 等价于 \(Q'' \tau N'\)；于是 \(Q \tau (M \tau N)\) 等价于 \((Q \tau M) \tau N\)，从而等价于 \((Q' \tau M') \tau N\)，也等价于 \((Q' \tau N) \tau M'\)，从而又等价于 \((Q'' \tau N') \tau M'\) 与 \(Q'' \tau (N' \tau M')\)；由于 \(M \tau N\) 与 \(N' \tau M'\) 是中性的，确实有 \(Q \sim Q''\)。因此，关系 \(Q \sim Q'\) 是 \(Q\) 与 \(Q'\) 之间的一个等价关系。显然，若 \(Q\) 与 \(Q'\) 是两个有限维且等价的非退化二次型，则有 \(Q \sim Q'\)。

  对 A 上每个非退化且有限维的二次型 Q，我们置

\[
\theta (Q) = \tau_ {x} (X \sim Q),\tag{1}
\]

并称 \(\theta(Q)\) 为 \(Q\) 的型。若 \(Q\) 与 \(Q'\) 是 \(A\) 上两个非退化且有限维的二次型，则关系 \(Q \sim Q'\) 与 \(\theta(Q) = \theta(Q')\) 等价。

命题 1. --- 设 Q 与 Q' 是 A 上两个非退化、有限维的二次型。要使 Q 与 Q' 等价，必须且只需它们具有相同的维数与相同的型。

  这个条件显然是必要的。设它被满足。于是存在中性型 N、N′，使得 Q ⊥ N 与 Q′ ⊥ N′ 等价。由于这两个型有相同的维数，N 与 N′ 亦然，从而它们等价（§ 4 第 2 小节命题 2 的推论 2）。于是由维特定理，Q 与 Q′ 等价（§ 4 第 3 小节定理 1 的推论 1）。

命题 2. --- 关系「存在 A 上一个非退化、有限维的二次型 Q，使得 X = θ(Q)」在 X 中是集合化的（《集合论》第二章 § 1 第 4 小节）。

  事实上，设 V 是 A 上一个无限维向量空间，S 是定义在 V 的有限维子空间上的非退化二次型所成的集合，W 是诸 \(\theta(Q)\)（\(Q \in \mathfrak{S}\)）所成的集合。显然，A 上每个非退化、有限维的二次型 \(Q'\) 等价于 \(\mathfrak{S}\) 的至少一个元素；由此 \(\theta(Q') \in \mathfrak{W}\)，这证明了我们的断言。

\subsection{二次型群.}

  我们将给 A 上有限维非退化二次型的型所成的集合 \(\mathfrak{W}\) 装备一个交换群结构。我们将用公式在 \(\mathfrak{W}\) 中定义一个加法

\[
\mathrm{T} + \mathrm{T} ^ {\prime} = \theta (\mathrm{T} _ {\tau} \mathrm{T} ^ {\prime})\tag{2}
\]

\[
(\mathrm{T}, \mathrm{T} ^ {\prime} \text {   in   } \mathfrak {W}),
\]

  这个加法是交换的，因为 \(T^{\prime} \tau T\) 等价于 \(T \tau T^{\prime}\)。它是结合的，因为若 T、\(T^{\prime}\) 与 \(T^{\prime\prime}\) 是 W 的元素，我们有

\[
\begin{array}{r l} & (\mathrm{T} + \mathrm{T} ^ {\prime}) + \mathrm{T} ^ {\prime \prime} \sim (\mathrm{T} + \mathrm{T} ^ {\prime}) \tau \mathrm{T} ^ {\prime \prime} \sim (\mathrm{T} \tau \mathrm{T} ^ {\prime}) \tau \mathrm{T} ^ {\prime \prime} \\ & \qquad \sim \mathrm{T} \tau (\mathrm{T} ^ {\prime} \tau \mathrm{T} ^ {\prime \prime}) \sim \mathrm{T} \tau (\mathrm{T} ^ {\prime} + \mathrm{T} ^ {\prime \prime}) \sim \mathrm{T} + (\mathrm{T} ^ {\prime} + \mathrm{T} ^ {\prime \prime}), \end{array}
\]

从而 \((\mathrm{T} + \mathrm{T}' ) + \mathrm{T}'' = \mathrm{T} + (\mathrm{T}' + \mathrm{T}'' )\)，因为 W 中具有相同型的两个元素相等。此外，刚才定义的加法有一个单位元素：事实上，显然所有中性型都有同一个型 \(T_{0}\)，即维数为零的零型的型；立即可见 \(T_{0}\) 就是所求的中性元素。最后，对每个 \(T \in W\) 都存在与 T 相反的元素，这一点立即由下列命题得出：

命题 3. — 设 Q 是 A 上向量空间 V 上一个非退化、有限维的二次型。用 −Q 记 V 上由 \((-Q)(x) = -Q(x) (x \in V)\) 定义的二次型。则型 \(Q \tau (-Q)\) 是中性的。

  事实上，\(Q \tau (-Q)\) 在对角线 \(D\)（\(V \times V\) 的对角线）上的限制是零。因此这个型的指标 \(\geqslant\frac{1}{2}\dim(V\times V)\)（ \(\S\) 4 第 2 小节定义 2），从而等于 \(\frac{1}{2}\dim(V\times V)\)（同出处公式 (4)）。由此得出 \(Q\tau(-Q)\) 是中性的（同出处）。

  这使我们能够陈述下列定义：

定义 1. --- A 上有限维非退化二次型的型所成的集合，带上由 (2) 定义的加法，称为 A 的二次型群，或 A 的维特群。

注. --- 1) 每个型为零（即用上面的记号使 \(\theta(Q) = T_0\)）的非退化、有限维二次型 Q 都是中性型。事实上，存在中性型 N、N′，使得 \(Q \tau N\) 等价于 N′。这表明 Q 是偶维的，从而存在一个与 Q 同维数的中性型 \(N_1\)。由于 Q 与 \(N_1\) 有相同的型，由命题 1 它们等价，故 Q 是中性的。

\begin{enumerate}
\def\labelenumi{\arabic{enumi})}
\setcounter{enumi}{1}
\tightlist
\item
  对 A 上每个有限维二次型 Q，用 \(\delta(Q)\) 记 Q 的维数模 2 的类。我们有
\end{enumerate}

\[
\delta (Q \tau Q ^ {\prime}) = \delta (Q) + \delta (Q ^ {\prime}).
\]

  由于每个中性型 N 都是偶维的，我们有 \(\delta(N) = 0\)；因此关系 \(Q \sim Q'\) 蕴涵 \(\delta(Q) = \delta(Q')\)。于是 \(\delta\) 在 A 上二次型的型所成的群 \(\mathfrak{W}\) 上的限制是 \(\mathfrak{W}\) 到群 \(\mathbf{Z}/(2)\) 中的一个同态。当 A 的特征 \(\neq 2\) 时这个同态是满射的，但当 A 的特征为 2 时则不然，因为这时奇维的二次型是退化的，由于其相伴双线性型是交错的（参见 § 5）。

\begin{enumerate}
\def\labelenumi{\arabic{enumi})}
\setcounter{enumi}{2}
\tightlist
\item
  设 \(a\) 是 A 的一个 \(\neq 0\) 的元素。若 N 是中性型，则 \(aN\) 也是。由此得出关系 \(Q \sim Q'\) 蕴涵 \(aQ \sim aQ'\)。对群 \(\mathfrak{W}\) 的每个元素 T，我们置
\end{enumerate}

\[
a. \mathrm{T} = \theta (a \mathrm{T}).\tag{3}
\]

  这样我们就在 A 的非零元素所成的群 A* 与群 W 之间得到一个外合成法则。下列公式立即由定义得出：

\[
a. (\mathrm{T} + \mathrm{T} ^ {\prime}) = a. \mathrm{T} + a. \mathrm{T} ^ {\prime}, \quad a b. \mathrm{T} = a. (b. \mathrm{T})\tag{4}
\]

\[
(a, b \text {   in   } A ^ {*}, T, T ^ {\prime} \text {   in   } \mathfrak {W}).
\]

  另一方面，若 \(a\)、\(b\) 与 \(a + b\) 都属于 \(\mathbf{A}^*\)，则一般没有 \((a + b)\) . \(\mathrm{T} = a\) . \(\mathrm{T} + b\) . \(\mathrm{T}\)（ \(\mathrm{T} \in \mathfrak{W}\)）。

命题 4. --- 设 Q 是 A 上有限维向量空间 E 上一个非退化二次型。假设 A 的特征 \(\neq 2\)，且设 \((x_{1},\ldots,x_{n})\) 是 V 的一个正交基。用 \(T_{1}\) 记二次型 \(Q_{1}\)（定义在向量空间 A 上且使 \(Q_{1}(1)=1\)）的型。则 Q 的型是 \(\sum_{i=1}^{n}\mathrm{Q}(x_{i}).T_{1}\)。

  事实上，型 Q 等价于

\[
(\mathrm{Q} (x _ {1}) \mathrm{Q} _ {1}) \tau \dots \tau (\mathrm{Q} (x _ {n}) \mathrm{Q} _ {1})
\]

推论. --- 在命题 3 的假设与记号下，元素 \(a\) . \(\mathrm{T}_{1}\)（\(a \in \mathrm{A}^{*}\)）构成 A 上二次型群的一组生成元。

  这样，确定 A 上二次型群的结构就归结为求形如 \(a.T_{1}\) 的元素之间存在的 Z-线性关系。若 \(b \in A^{*}\)，命题 4 中定义的型 \(Q_{1}\) 显然等价于 \(b^{2}Q_{1}\)；因此有 \(a.T_{1} = ab^{2}.T_{1}\)，这表明 \(a.T_{1}\) 只依赖于 a 模子群 \((A^{*})^{2}\)（由 \(A^{*}\) 中元素的平方组成）的类。此外，由命题 3 得出 \((-a).T_{1} = -a.T_{1}\)。然而，除了由刚才指出的关系推出的那些之外，诸 \(a.T_{1}\) 之间一般还存在别的 Z-线性关系。

命题 5. --- 假设 A 是一个极大有序域。设 Q 是 A 上一个非退化、有限维的二次型，(s, t) 是它的符号差（§ 7 第 2 小节定义 2）。则 Q 的型是 (s - t)·T₁，且 A 上二次型的群 W 是由 T₁ 生成的无限循环群。

  事实上，由于 \(\mathbf{A}^{*}/(\mathbf{A}^{*})^{2}\) 的阶为 2，且 \((-1)\) . \(\mathrm{T}_{1} = -\mathrm{T}_{1}\)，故 \(\mathfrak{W}\) 由 \(\mathrm{T}_{1}\) 生成，从而是单演的。对一切 \(n > 0\)，\(n\) . \(\mathrm{T}_{1}\) 是维数 \(n\) 的正非退化二次型的型；由于这些型不是中性的，我们有 \(n\) . \(\mathrm{T}_{1} \neq 0\)，这表明 \(\mathfrak{W}\) 是无限的。最后，用命题 4 的记号，一个符号差为 \((s, t)\) 的型同构于 s 个 \(Q_{1}\) 与 t 个 \(-Q_{1}\) 的直和（§ 7 第 2 小节定理 1）；由此得出它的型是 \((s - t)\) . \(T_{1}\)。

\subsection{二次型环.}

在本小节中我们将假设 A 是一个特征 ≠ 2 的体。

  给定 A 上向量空间 V、V' 上的两个二次型 Q、Q'，我们称 V ⊗ V' 上的、其相伴双线性型是与 Q、Q' 相伴的双线性型的张量积（§ 1 第 9 小节定义 11）的那个二次型为 Q 与 Q' 的张量积，并用 Q ⊗ Q' 记之。容易看出 Q ⊗ Q' 满足关系

\[
(Q \otimes Q ^ {\prime}) (x \otimes x ^ {\prime}) = Q (x) Q ^ {\prime} (x ^ {\prime}) \quad (x \in V, x ^ {\prime} \in V ^ {\prime}).\tag{5}
\]

  若 Q 与 Q' 都是非退化且有限维的，则 Q ⊗ Q' 亦然（§ 1 第 9 小节命题 9）。

  设 Q、Q′、Q″ 是向量空间 V、V′、V″ 上的二次型。利用 V ⊗ V' 到 V' ⊗ V（相应地，(V ⊗ V′) ⊗ V″ 到 V ⊗ (V′ ⊗ V″)、(V × V′) ⊗ V″ 到 (V ⊗ V″) × (V′ ⊗ V″)）上的典范同构，我们立即看出 Q ⊗ Q' 等价于 Q' ⊗ Q（相应地，(Q ⊗ Q′) ⊗ Q″ 等价于 Q ⊗ (Q′ ⊗ Q″)、(Q τ Q′) ⊗ Q″ 等价于 (Q ⊗ Q″) τ (Q′ ⊗ Q″)）。

  设 Q 与 Q' 是两个非退化、有限维的二次型。若 Q 是中性的，则 Q⊗Q' 亦然。事实上，设 V、V' 是 Q、Q' 的定义空间，维数分别为 2n 与 n'，并设 W 是 V 的一个维数为 n 的全奇异子空间（§ 4 第 2 小节）；则 W⊗V' 是一个全奇异子空间，其维数是 V⊗V' 的维数的一半；由此，如同命题 3 那样推出 Q⊗Q' 是中性的。同样，只要 Q' 是中性的，Q⊗Q' 就是中性的。

  由此我们推出：若 Q、Q'、Q₁、Q₁' 是 A 上非退化、有限维的二次型，且假设 θ(Q₁) = θ(Q) 与 θ(Q₁') = θ(Q')，则有 θ(Q₁ ⊗ Q₁') = θ(Q ⊗ Q')。事实上，只需在 \(Q_{1} = Q \tau N\) 与 \(Q_{1}' = Q' \tau N'\)（N 与 \(N'\) 是中性型）的情形验证这一点；这时 \(Q_{1} \otimes Q_{1}'\) 等价于

\[
(Q \otimes Q ^ {\prime}) \tau (Q \otimes N ^ {\prime} \tau Q ^ {\prime} \otimes N \tau N \otimes N ^ {\prime})
\]

而第二个括号表示一个中性型；这就证明了我们的断言。

  现在设 \(\mathfrak{W}\) 是 A 上二次型的型所成的群。让我们在集合 \(\mathfrak{W}\) 上用公式定义第二个合成法则，它用乘法记号表示

\[
\mathrm{TT} ^ {\prime} = \theta (\mathrm{T} \otimes \mathrm{T} ^ {\prime})\tag{6}
\]

\[
(\mathbf {T}, \mathbf {T} ^ {\prime} \text {   in   } \mathfrak {W}).
\]

  由刚才所见立即得出，这个合成法则是交换的、结合的，且关于加法是分配的。它有一个单位元素，即型 \(T_{1}\)（二次型 \(Q_{1}\) 的型，该二次型定义在向量空间 A 上且使 \(Q_{1}(1)=1\)）：事实上，由 (5)，对每个二次型 Q 有 \(Q_{1}\otimes Q=Q\)。因此，加法群 W 带上刚才定义的乘法是一个有单位元素的交换环；它称为 A 的二次型环（或 A 的维特环，当不致混淆时）。

注. --- 1) 显然，若 a 是 A* 的一个元素，则

\[
a. (\mathrm{TT} ^ {\prime}) = (a. \mathrm{T}) \mathrm{T} ^ {\prime} = \mathrm{T} (a. \mathrm{T} ^ {\prime})\tag{7}
\]

对 W 的任意元素 T、T' 成立。此外应注意有 a.T = TaT，其中用 Ta 记 A 上二次型 aQ1 的型。

\begin{enumerate}
\def\labelenumi{\arabic{enumi})}
\setcounter{enumi}{1}
\tightlist
\item
  由于 A 的特征 \(\neq 2\)，\(\mathfrak{W}\) 的每个元素 T 都有 \(\sum_{i=1}^{n} a_i \cdot T_1\) 的形式，其中 \(a_i \in A^*\)（第 2 小节命题 4）。我们有
\end{enumerate}

\[
(\sum_ {i = 1} ^ {n} a _ {i}. \mathrm{T} _ {1}) (\sum_ {j = 1} ^ {q} b _ {j}. \mathrm{T} _ {1}) = \sum_ {i, j} a _ {i} b _ {j}. \mathrm{T} _ {1} \quad (a _ {i}, b _ {j} \text {   in   } \mathrm{A} ^ {*}).\tag{8}
\]

\begin{enumerate}
\def\labelenumi{\arabic{enumi})}
\setcounter{enumi}{2}
\tightlist
\item
  假设 A 是一个极大有序域。则环 W 同构于 Z（命题 5），与符号差为 \((s, t)\) 的型的型对应的整数是 s - t（同出处）。由于两个型 Q、\(Q'\)（符号差分别为 \((s, t)\)、\((s', t')\)）的张量积是维数为 \((s + t)(s' + t')\) 的型，故由一个初等计算得出 \(Q \otimes Q'\) 的符号差是 \((ss' + tt', st' + ts')\)。
\end{enumerate}

\section{克利福德代数}

在本节中我们将假设环 A 是交换的。我们将用 Q 记 A-模 E 上的一个二次型，用 \(\Phi\) 记相伴的双线性型（ \(\S\) 3 第 4 小节）。

\subsection{克利福德代数的定义与泛性质.}

定义 1. --- 我们称模 E 的张量代数 T(E) 对由形如 \(x \otimes x - Q(x)\) . 1（\(x \in E\)）的元素生成的双边理想（记作 I(Q)）的商代数为 Q 的克利福德代数，并用 C(Q) 记之。

  我们将用 \(\rho_{\mathbb{Q}}\)（当不致混淆时也简记作 \(\rho\)）记由 E 到 T(E) 中的典范映射与 T(E) 到 C(Q) 上的典范映射 \(\sigma\) 复合而得的从 E 到 C(Q) 中的映射；映射 \(\rho_{\mathbb{Q}}\) 称为典范的。注意 C(Q) 由 \(\rho_{\mathbb{Q}}(E)\) 生成，且对 \(x \in E\)，有

\[
\rho (x) ^ {2} = \mathrm{Q} (x). 1;\tag{1}
\]

从而，把 \(x\) 换成 \(x + y\)（ \(x, y\) 属于 E）

\[
\rho (x) \rho (y) + \rho (y) \rho (x) = \Phi (x, y). 1\tag{2}
\]

例. --- 若 E 承认一个由单个元素 \(e\) 组成的基，则 T(E) 同构于多项式代数 A{[}X{]}，而 C(Q) 是 A 的一个二次扩张，以 \((1, u)\) 为基，其中 \(u\) 是元素 \(u = \rho(e)\)，它满足 \(u^2 = Q(e)\)。

  用 \(T^{h}\) 记 h 次张量幂 \(\otimes E\)（在 \(\mathrm{T}(\mathrm{E})\) 中），并设 \(T^{+}\)（相应地，\(T^{-}\)）是 h 为偶数（相应地，奇数）的那些 \(T^{h}\) 之和。

  由于 T(E) 是 \(T^{+}\) 与 \(T^{-}\) 的直和，且 I(Q) 由 \(T^{+}\) 的元素生成，故 I(Q) 是 \(T^{+} \cap I(Q)\) 与 \(T^{-} \cap I(Q)\) 的直和，而 C(Q) 是两个子模 \(C^{+}(Q) = \sigma(T^{+})\) 与 \(C^{-}(Q) = \sigma(T^{-})\)（也记作 \(C^{+}\) 与 \(C^{-}\)）的直和。\(C^{+}\) 的元素将称为偶的（相应地，奇的）。我们有关系式

\[
\mathrm{C} ^ {+} \mathrm{C} ^ {+} \subset \mathrm{C} ^ {+}, \quad \mathrm{C} ^ {+} \mathrm{C} ^ {-} \subset \mathrm{C} ^ {-}, \quad \mathrm{C} ^ {-} \mathrm{C} ^ {+} \subset \mathrm{C} ^ {-}, \quad \mathrm{C} ^ {-} \mathrm{C} ^ {-} \subset \mathrm{C} ^ {+}.\tag{3}
\]

  特别地，\(C^{+}\) 是 \(C(Q)\) 的一个子代数。

命题 1. --- 设 f 是 E 到 A 上一个代数 D 中的线性映射，使得 \(f(x)^{2} = \mathrm{Q}(x)\) . 1 对每个 \(x \in \mathbf{E}\) 成立。存在一个且只有一个 C(Q) 到 D 中的同态 \(\bar{f}\)，使得 \(f = \bar{f} \circ \rho_{\mathbb{Q}}\)。

  \(\bar{f}\) 的唯一性由 C(Q) 由 \(\rho_{\mathbb{Q}}(\mathrm{E})\) 生成这一事实得出。设 \(h\) 是扩张 \(f\) 的 T(E) 到 D 中的唯一同态（\(h\) 由 \(h(x_1 \otimes \cdots \otimes x_n) = f(x_1) \ldots f(x_n)\) 定义）。我们有

\[
h (x \otimes x - \mathrm{Q} (x). 1) = (f (x) ^ {2} - \mathrm{Q} (x)). 1 = 0,
\]

  因而 h 在 I(Q) 上为零，并通过过渡到商定义所求的同态 \(\bar{f}\)。

  命题 1 表明 C(Q) 是一个泛映射问题的解（《集合论》第四章 § 3 第 1 小节）。

  特别地，取 D 为 C(Q) 的反代数，取 \(f\) 为映射 \(\rho\)；命题 1 蕴涵存在一个且只有一个 C(Q) 的反自同构 \(\beta\)，它在 \(\rho(E)\) 上的限制是恒等映射；它称为 C(Q) 的主反自同构。显然 \(\beta^{2} = 1\)。

  另一方面，设 Q' 是 A-模 E' 上的一个二次型，f 是 E 到 E' 中的线性映射，使得 Q'∘f = Q。我们有 \(\rho_{\mathbb{Q}'}(f(x))^{2} = \mathrm{Q}'(f(x)).1 = \mathrm{Q}(x).1\)，因而存在一个且只有一个从 C(Q) 到 C(Q') 中的同态 C(f)，使得 C(f)∘ρq = ρq'∘f。若 f 是恒等映射，则 C(f) 是恒等映射；若 Q'' 是 A-模 E'' 上的一个二次型，g 是 E' 到 E'' 中的线性映射，使得 Q''∘g = Q'，则我们有 C(g∘f) = C(g)∘C(f)。当 E' 是 E 的子模且 f 是 E' 到 E 中的典范单射（于是 Q' 是 Q 在 E' 上的限制）时，我们称 C(f) 为 C(Q') 到 C(Q) 中的典范同态。

  特别地，取 Q' = Q 并取 f 为映射 \(x \to -x\)；我们看出存在一个且只有一个 C(Q) 的自同构 \(\alpha\)，使得 \(\alpha \circ \rho = -\rho\)；它称为 C(Q) 的主自同构。显然 \(\alpha^{2} = 1\)，且 \(\alpha\) 在 C⁺（相应地，C⁻）上的限制是恒等映射（相应地，映射 \(u \to -u\)）。

命题 2. --- 设 A' 是一个交换环，φ 是 A 到 A' 中的同态，Q' 是由 Q 经纯量扩张得到的 E' = A' ⊗\_A E 上的二次型（§ 3 第 4 小节命题 3）。存在一个且只有一个代数 A' ⊗\_A C(Q) 到 C(Q') 上的同构 j，使得对一切 x ∈ E 有 j (1 ⊗ ρ\_Q(x)) = ρ\_Q′(1 ⊗ x)。

  只需证明代数 \(C' = A' \otimes C(Q)\) 与映射 \(1 \otimes_{\rho_{Q}}\)（从 \(E'\) 到 \(C'\) 中）构成与 \(C(Q')\) 和 \(\rho_{Q'}\) 相同的泛问题的一个解。事实上，设 \(D'\) 是 \(A'\) 上一个代数，\(f'\) 是 \(E'\) 到 \(D'\) 中的一个 A'-线性映射，使得 \(f'(x')^{2} = Q'(x')\) . 1 对每个 \(x' \in E'\) 成立。映射 \(g : x \to f'(1 \otimes x)\)（从 E 到 \(D'\) 中，后者经由同态 \(\varphi\) 看作 A-模）是 A-线性的，且 \(g(x)^{2} = Q'(1 \otimes x)\) . 1 = Q(x). 1 对一切 \(x \in E\) 成立。因此存在一个且只有一个 A-同态 \(\bar{g}\)（从 C(Q) 到 \(D'\) 中），使得 \(\bar{g}(\rho_{Q}(x)) = f'(1 \otimes x)\)。从而存在一个且只有一个 A'-同态 \(\overline{f'}\)（从 C' 到 \(D'\) 中），使得 \(\overline{f}'(1 \otimes \rho_{Q}(x)) = f'(1 \otimes x)\) 对一切 \(x \in E\) 成立；由线性性得出，\(\overline{f}'((1 \otimes \rho_{Q})(x')) = f'(x')\) 对一切 \(x' \in E'\) 成立。证毕。

\subsection{张量代数中的某些运算.}

在本小节中，我们将用 \(e_x\)（\(x \in E\)）记张量代数 T(E) 到自身的线性映射 \(u \to x \otimes u\)。

引理 1. — 设 f 是 E 的对偶 E* 的一个元素。存在一个且只有一个 T(E) 到自身的线性映射 \(i_{f}\)，使得

\begin{enumerate}
\def\labelenumi{(\arabic{enumi})}
\setcounter{enumi}{3}
\tightlist
\item
\end{enumerate}

\[
i _ {f} (1) = 0\tag{5}
\]

\[
i _ {f} \circ e _ {x} + e _ {x} \circ i _ {f} = f (x). I
\]

对一切 \(x \in \mathbf{E}\)

（其中 I 表示恒等映射）。映射 \(f \to i_f\)（从 E* 到 \(\mathfrak{U}(\mathrm{T}(\mathrm{E}))\) 中）是线性的。有 \(i_f(\mathrm{T}^n) \subset \mathrm{T}^{n-1}\)，\((i_f)^2 = 0\)，且 \(i_f \circ i_g + i_g \circ i_f = 0\) 对 E* 中的 \(f\)、\(g\) 成立。映射 \(i_f\) 在子代数

T(E)（由 f 的核生成的）上为零。理想 I(Q) 在 \(i_{f}\) 下稳定；因此通过过渡到商，\(i_{f}\) 定义 C(Q) 到自身的一个线性映射（仍记作 \(i_{f}\)）。

  事实上，公式 (5) 可以写成

\[
i _ {f} (x \otimes u) = - x \otimes i _ {f} (u) + f (x). u \quad (x \in \mathrm{E}, u \in \mathrm{T(E)}).\tag{6}
\]

  由于 (4) 完全确定了 \(i_{f}\) 在 \(\mathbf{T}^{0}\) 上的值，且 (6) 确定了 \(i_{f}\) 在 \(\mathbf{T}^{n}\) 上的值（当其在 \(\mathbf{T}^{n-1}\) 上的值已知时），\(i_{f}\) 的唯一性得证。另一方面，对 \(x\in\mathbf{E}\) 与 \(u\in\mathbf{T}^{n-1}\)，(6) 的右边在 \(\mathbf{E}\times\mathbf{T}^{n-1}\) 上是双线性的；这证明了 \(i_{f}\) 的存在性（对 \(n\) 作归纳；第三章 § 1 第 2 小节），并且对 \(n\) 作归纳可知 \(i_{f}(\mathbf{T}^{n})\subset\mathbf{T}^{n-1}\)。若 \(f=ag+bh\)（\(a,b\) 属于 A，\(g,h\) 属于 E*），则显然 \(ai_{g}+bi_{h}\) 满足 (4) 与 (5)，从而等于 \(i_{f}\)。我们有

\[
(i _ {f}) ^ {2} \circ e _ {x} = - i _ {f} \circ e _ {x} \circ i _ {f} + f (x) i _ {f} = e _ {x} \circ (i _ {f}) ^ {2}
\]

  又，由于 \((i_f)^2(1) = 0\)，我们对 \(n\) 作归纳推出 \((i_f)^2\) 在 \(\mathbf{T}^n\) 上为零。把 \(f\) 换成 \(f + g\)，关系 \((i_f)^2 = 0\) 就给出 \(i_f \circ i_g + i_g \circ i_f = 0\)。一个与前面类似的归纳论证表明 \(i_f\) 在由 \(f\) 的核生成的子代数上为零。最后，(6) 蕴涵 I(Q) 中元素 \(u\)（使 \(i_f(u) \in \mathrm{I}(Q)\)）所成之集是 T(E) 的一个左理想；此外，若 \(u = (x \otimes x - Q(x).1) \otimes v (x \in E, v \in T(E))\)，则有

\[
\begin{array}{r l} i _ {f} (u) & = f (x) x \otimes v - x \otimes i _ {f} (x \otimes v) - \mathrm{Q} (x) i _ {f} (v) \\ & = f (x) x \otimes v - f (x) x \otimes v + x \otimes x \otimes i _ {f} (v) - \mathrm{Q} (x) i _ {f} (v) \\ & = (x \otimes x - \mathrm{Q} (x). 1) \otimes i _ {f} (v); \end{array}
\]

  因此 I(Q) 在 \(i_{f}\) 下稳定，最后的断言随之得出。证毕。

  设 F 是 E 上的一个双线性型；在本节余下部分，我们将用 \(i_x^{\mathrm{F}}\)（\(x \in \mathrm{E}\)）记映射 \(i_f\)，它对应于 E 上的线性型 \(f: y \to \mathrm{F}(x, y)\)。

引理 2. --- 存在 T(E) 到自身的一个线性映射 \(\lambda_{\mathrm{F}}\)，使得

\begin{enumerate}
\def\labelenumi{(\arabic{enumi})}
\setcounter{enumi}{6}
\tightlist
\item
\end{enumerate}

\[
\lambda_ {\mathrm{F}} (1) = 1\tag{8}
\]

\[
\lambda_ {\mathrm{F}} \circ e _ {x} = (e _ {x} + i _ {x} ^ {\mathrm{F}}) \circ \lambda_ {\mathrm{F}}\tag{\( (x \in E) \) .}
\]

  无论 \(f \in \mathbf{E}^*\) 如何，都有

\[
\lambda_ {\mathrm{F}} \circ i _ {f} = i _ {f} \circ \lambda_ {\mathrm{F}}.\tag{9}
\]

  事实上，公式 (8) 等价于

\[
\lambda_ {\mathbb {F}} (x \otimes u) = x \otimes \lambda_ {\mathbb {F}} (u) + i _ {x} ^ {\mathrm{F}} (\lambda_ {\mathbb {F}} (u)) \quad (x \in \mathrm{E}, u \in \mathrm{T} (\mathrm{E})).\tag{10}
\]

  由于 (7) 完全确定了 \(\lambda_{\mathbb{F}}\) 在 \(\mathrm{T}^{0}\) 上的值，且 (10) 确定了 \(\lambda_{\mathbb{F}}\) 在 \(\mathrm{T}^{n}\) 上的值（当其在 \(\mathrm{T}^{n-1}\) 上的值已知后），\(\lambda_{\mathbb{F}}\) 的唯一性得证。另一方面，对 \(x \in \mathrm{E}\) 与 \(u \in \mathrm{T}^{n-1}\)，(10) 的右边在 \(\mathrm{E} \times \mathrm{T}^{n-1}\) 上是双线性的；这证明了 \(\lambda_{\mathbb{F}}\) 的存在性（对 \(n\) 作归纳）。剩下的要证明 (9)，我们将用归纳法来做；(9) 的两边在 \(\mathrm{T}^{0}\) 上为零；设 (9) 在 \(\sum_{h=0}^{n-1} \mathrm{T}^{h}\) 上成立；于是对 \(x \in \mathrm{E}\) 与 \(u \in \mathrm{T}^{n-1}\) 我们有：

\[
\begin{array}{r l} (\lambda_ {\mathbb {F}} \circ i _ {f}) (x \otimes u) & = (- \lambda_ {\mathbb {F}} \circ e _ {x} \circ i _ {f} + f (x) \lambda_ {\mathbb {F}}) (u) \\ & = - (e _ {x} + i _ {x} ^ {\mathrm{F}}) \circ \lambda_ {\mathbb {F}} \circ i _ {f} (u) + f (x) \lambda_ {\mathbb {F}} (u) \\ & = - (e _ {x} + i _ {x} ^ {\mathrm{F}}) \circ i _ {f} \circ \lambda_ {\mathbb {F}} (u) + f (x) \lambda_ {\mathbb {F}} (u) \\ & = (i _ {f} \circ e _ {x} \circ \lambda_ {\mathbb {F}} - f (x) \lambda_ {\mathbb {F}} + i _ {f} \circ i _ {x} ^ {\mathrm{F}} \circ \lambda_ {\mathbb {F}} + f (x) \lambda_ {\mathbb {F}}) (u) \\ & = (i _ {f} \circ (e _ {x} + i _ {x} ^ {\mathrm{F}}) \circ \lambda_ {\mathbb {F}}) (u) = (i _ {f} \circ \lambda_ {\mathbb {F}}) (x \otimes u), \end{array}
\]

  从而得到我们的最后一个断言。

引理 3. --- 设 F 与 G 是 E 上两个双线性型。则我们有 \(\lambda_{\mathrm{F}} \circ \lambda_{\mathrm{G}} = \lambda_{\mathrm{F+G}}\)。对 E 上每个双线性型 F，\(\lambda_{\mathrm{F}}\) 是 T(E) 到自身上的一个双射。

  事实上，\(\lambda_{\mathbb{F}} \circ \lambda_{\mathbb{G}}\) 具有 \(\lambda_{\mathbb{F} + \mathbb{G}}: \text{有 } (\lambda_{\mathbb{F}} \circ \lambda_{\mathbb{G}})(1) = 1\) 的特征性质 (7) 与 (8)，并且

\[
\begin{array}{r l} \lambda_ {\mathrm{F}} \circ \lambda_ {\mathrm{G}} \circ e _ {x} & = \lambda_ {\mathrm{F}} \circ (e _ {x} + i _ {x} ^ {\mathrm{G}}) \circ \lambda_ {\mathrm{G}} = (e _ {x} + i _ {x} ^ {\mathrm{G}} + i _ {x} ^ {\mathrm{F}}) \circ \lambda_ {\mathrm{F}} \circ \lambda_ {\mathrm{G}} \\ & = (e _ {x} + i _ {x} ^ {\mathrm{F+G}}) \circ \lambda_ {\mathrm{F}} \circ \lambda_ {\mathrm{G}}. \end{array}
\]

  另一方面，若 F = 0，则 \(i_x^{\mathbb{F}} = 0\)（对一切 \(x \in \mathbb{E}\)），从而 \(\lambda_{\mathbb{F}} = \mathbb{I}\)，这蕴涵对每个 F 有 \(\lambda_{\mathbb{P}} \circ \lambda_{-\mathbb{P}} = \lambda_{-\mathbb{F}} \circ \lambda_{\mathbb{F}} = \mathbb{I}\)。

\subsection{克利福德代数的基.}

命题 3. --- 设 Q 与 Q' 是两个二次型，F 是 E 上一个双线性型，使得对一切 x ∈ E 有 Q'(x) = Q(x) + F(x, x)。映射 λ\_F 把理想 I(Q') 映到理想 I(Q) 上，并定义 A-模 C(Q') 到 A-模 C(Q) 上的一个同构（记作 \(\overline{\lambda_{\mathbf{F}}}\)）。

  由于 \(\lambda_{\mathbb{F}}\) 是一个双射，其逆双射是 \(\lambda_{-\mathbb{F}}\)（引理 3），只需证明包含关系 \(\lambda_{\mathbb{F}}(\mathrm{I}(\mathrm{Q}^{\prime})) \subset \mathrm{I}(\mathrm{Q})\)。由于 \(\mathrm{I}(\mathrm{Q})\) 是在 \(i_x^{\mathbb{F}}\) 下稳定的左理想（引理 1），(8) 表明由 \(u \in \mathrm{T}(\mathrm{E})\)（使 \(\lambda_{\mathbb{F}}(u) \in \mathrm{I}(\mathrm{Q})\)）组成的集是一个左理想。因此只需证明：对任意 \(u \in \mathrm{T}(\mathrm{E})\) 与 \(x \in \mathrm{E}\)，有 \(\lambda_{\mathbb{F}}(x \otimes x \otimes u - \mathrm{Q}'(x)u) \in \mathrm{I}(\mathrm{Q})\)。现在，由 (8) 与引理 1，我们有

\[
\lambda_ {\mathrm{F}} \circ e _ {x} ^ {2} = (e _ {x} + i _ {x} ^ {\mathrm{F}}) ^ {2} \circ \lambda_ {\mathrm{F}} = (e _ {x} ^ {2} + \mathrm{F} (x, x)) \circ \lambda_ {\mathrm{F}},
\]

从而

\[
\begin{array}{r l} \lambda_ {\mathbf {F}} (x \otimes x \otimes u - \mathrm{Q} ^ {\prime} (x) u) & = (e _ {x} ^ {2} + \mathrm{F} (x, x) - \mathrm{Q} ^ {\prime} (x)) \circ \lambda_ {\mathbf {F}} (u) \\ & = (x \otimes x - \mathrm{Q} (x)) \otimes \lambda_ {\mathbf {F}} (u) \in \mathrm{I} (\mathrm{Q}). \end{array}
\]

引理 4. — 若二次型 Q 为零，则 C(Q) 不是别的，正是 E 的外代数。

  事实上，E 的外代数不是别的，正是 T(E) 对由形如 \(a \otimes v \otimes a\)（\(a \in E\)，\(v \in T(E)\)）的元素生成的双边理想 J 的商（第三章 § 5 第 5 小节与第 9 小节）。显然 I(Q) ⊂ J。因此只需证明 \(a \otimes v \otimes a \in I(Q)\)。若 \(v \in T^0\)，这是明显的。设这个断言对 \(n-1 \leqslant \nu \in \sum_{h=0}^{\infty} \mathrm{T}^{h}\) 已被证明，并设 \(x \in \mathrm{E}\) 与 \(u \in \mathrm{T}^{n-1}\)；我们有

\[
\begin{array}{c} \boldsymbol {a} \otimes \boldsymbol {x} \otimes \boldsymbol {u} \otimes \boldsymbol {a} = (\boldsymbol {a} + \boldsymbol {x}) \otimes (\boldsymbol {a} + \boldsymbol {x}) \otimes \boldsymbol {u} \otimes \boldsymbol {a} \\ - \boldsymbol {a} \otimes \boldsymbol {a} \otimes \boldsymbol {u} \otimes \boldsymbol {a} - \boldsymbol {x} \otimes \boldsymbol {a} \otimes \boldsymbol {u} \otimes \boldsymbol {a} - \boldsymbol {x} \otimes \boldsymbol {x} \otimes \boldsymbol {u} \otimes \boldsymbol {a} \end{array}
\]

  而右边的四项都属于 I(Q)。

  特别地，假设模 E 承认一个基 \((x_{i})_{i\in \mathbf{L}}\)，并让我们把指标集 L 全序化。我们知道（第三章 § 5 第 6 小节）E 的外代数承认由元素 \(x_{\mathrm{H}}\) 组成的族为基，其中 H 遍历 L 的有限子集所成之集，而 \(x_{\mathrm{H}}\) 是元素 \(x_{h_1}\wedge \ldots \wedge x_{h_q}\)，\((h_1,\ldots ,h_q)\) 表示 H 的元素的严格递增序列。

  另一方面，考虑由下列条件定义的双线性型 F：\(\mathrm{F}(x_{i}, x_{j}) = -\Phi(x_{i}, x_{j})\)（若 i \textgreater{} j），\(\mathrm{F}(x_{i}, x_{j}) = 0\)（若 i \textless{} j），而 \(\mathrm{F}(x_{i}, x_{i}) = -\mathrm{Q}(x_{i})\)。显然 \(\mathrm{Q}(x) + \mathrm{F}(x, x) = 0\)；用命题 3 的记号，由这个命题与引理 4 得出 \(\bar{\lambda}_{\mathrm{F}}\) 是 \(\wedge \mathrm{E}\) 到 C(Q) 上的一个同构，因而 C(Q) 是一个自由 A-模。让我们对 H 的元素个数作归纳来证明：我们有

\[
\overline {{{{\lambda}}}} _ {\mathrm{F}} (x _ {\mathrm{H}}) = \rho (x _ {h _ {1}}) \dots \rho (x _ {h _ {q}})\tag{11}
\]

\((\mathrm{where}\ (h_{1},\ldots,h_{q})\) 是 H 的元素的严格递增序列）。若 H 有 0 个或 1 个元素，这是明显的。设 (11) 对至多有 \(q-1\) 个元素的集合成立。考虑一个有 \(q\) 个元素的子集 H，用 \(j\) 记它的最小元素，并让我们写 \(\mathrm{H}=\{j\}\cup\mathrm{K}\)，其中 K 是有 \(q-1\) 个元素的子集。由 (8) 与归纳假设，我们有

\[
\overline {{{\lambda}}} _ {\mathrm{F}} (x _ {\mathrm{H}}) = \overline {{{\lambda}}} _ {\mathrm{F}} (x _ {j} \wedge x _ {\mathrm{K}}) = \rho (x _ {j}) \overline {{{\lambda}}} _ {\mathrm{F}} (x _ {\mathrm{K}}) + i _ {x _ {j}} ^ {\mathrm{F}} (\overline {{{\lambda}}} _ {\mathrm{F}} (x _ {\mathrm{K}})) = x _ {\mathrm{H}} ^ {\prime} + i _ {x _ {j}} ^ {\mathrm{F}} (x _ {\mathrm{K}} ^ {\prime}),
\]

这里对 L 的每个有限子集 J，置 \(x_{J}^{\prime}=\rho(x_{j_{1}})\ldots\rho(x_{j_{s}})\)，其中 \((j_{1},\ldots,j_{s})\) 是 J 的元素的递增序列。现在，对 \(i\in\mathbf{K}\)，\(x_{i}\) 属于线性型 \(y\to\mathrm{F}(x_{j},y)\) 的核，故 \(i_{x_{j}}^{\mathrm{F}}(x_{\mathrm{K}}^{\prime})=0\)（引理 1）。这就证明了所要的结果。

  这样我们就证明了下面的定理：

定理 1. --- 假设 A-模 E 承认一个基 \((x_{i})_{i\in L}\)，指标集 L 带有一个全序结构。对 L 的每个有限子集 H，置 \(x_{\mathrm{H}} = \rho(x_{h_1})\rho(x_{h_2})\ldots \rho(x_{h_q})\)，其中 \((h_1,\ldots,h_q)\) 是 H 的元素的严格递增序列。则诸元素 \(x_{\mathrm{H}}\) 构成 A-模 C(Q) 的一个基。

推论 1. --- 若 E 是维数为 n 的自由模，则 C(Q) 是维数为 \(2^{n}\) 的自由模；此外，若 \(n > 0\)，则 C \(^{+}\) 与 C \(^{-}\) 是维数为 \(2^{n-1}\) 的自由模。

  这由二项式系数的性质立即得出。

推论 2. --- 若 E 是自由模，则从 E 到 C(Q) 中的典范映射 \(\rho\) 与从 A 到 C(Q) 中的映射 \(a \to a\) .1 都是单射。

推论 3. --- 假设 E 是两个自由子模 \(E_{1}\) 与 \(E_{2}\) 的直和。设 \(Q_{i}\) 是 Q 在 \(E_{i}\) 上的限制，\(p_{i}\) 是从 C(Q \(_{i}\) ) 到 C(Q) 中的典范映射（i = 1, 2）。则

线性映射 \(p\)（从 \(\mathrm{C(Q_{1})\otimes C(Q_{2})}\) 到 \(\mathrm{C(Q)}\) 中）由双线性映射 \((a,b)\to p_{1}(a)p_{2}(b)\)（从 \(\mathrm{C(Q_{1})\times C(Q_{2})}\) 到 \(\mathrm{C(Q)}\) 中）导出，是一个双射。

  事实上，只需考虑由 \(E_{1}\) 的一个基与 \(E_{2}\) 的一个基取并所得到的 E 的基。

推论 4. --- 在推论 3 的假设与记号下，进一步假设 \(E_{1}\) 与 \(E_{2}\) 正交，并在 \(C(Q_{1})\otimes C(Q_{2})\) 上转移 \(C(Q)\) 的代数结构（借助双射 p）。若 \(a_{i}\) 与 \(b_{i}\) 是 \(C(Q_{i})\) 的偶元素或奇元素（i=1, 2），则我们有 \((a_{1}\otimes a_{2})(b_{1}\otimes b_{2})=\varepsilon(a_{1}b_{1})\otimes(a_{2}b_{2})\)，其中 \(\varepsilon=1\)，除非 \(a_{2}\) 与 \(b_{1}\) 都是奇的，此时 \(\varepsilon=-1\)。

  事实上，只需证明 \(p_2(a_2)p_1(b_1) = \varepsilon p_1(b_1)p_2(a_2)\)，为此可以假设 \(p_2(a_2)\)（相应地，\(p_1(b_1)\)）是元素的一个乘积 \(x_1 \ldots x_h\)（相应地，\(y_1 \ldots y_k\)），这些元素属于 \(\rho_{\mathbb{Q}}(\mathbf{E}_2)\)（相应地，\(\rho_{\mathbb{Q}}(\mathbf{E}_1)\)）。由于 \(\mathbf{E}_1\) 与 \(\mathbf{E}_2\) 正交，我们有

\[
x _ {i} y _ {j} + y _ {j} x _ {i} = \Phi (x _ {i}, y _ {j}) = 0,
\]

  从而

\[
x _ {1} \dots x _ {h} y _ {1} \dots y _ {k} = (- 1) ^ {h k} y _ {1} \dots y _ {k} x _ {1} \dots x _ {h}.
\]

  如果略去 \(\mathbf{E}_1\) 与 \(\mathbf{E}_2\) 是自由模这一假设，推论 3 与 4 的结论仍然成立（参见习题 1）。

\subsection{克利福德代数的结构.}

  在本节中，我们将假设 A 是一个域，E 是 A 上有限维 m 的向量空间，且二次型 Q 是非退化的（由 § 5 第 1 小节定理 1，当 A 的特征为 2 时这要求 m 为偶数）。由于 E 是自由的，典范映射 ρ 是单射（第 3 小节定理 1 的推论 2）。今后我们将把 E 与它在 C(Q) 中的象等同起来。

定理 2. --- 假设 E 的维数是偶数 \(m = 2r\)，且 Q 是中性的（\(\S 4, \text{n}^{\circ} 2\)）。则代数 C(Q) 是可分的（第八章 \(\S 7, \text{n}^{\circ} 5\)，定义 1），并且同构于 A 上维数为 \(2^{r}\) 的向量空间的自同态代数。此外，若 \(m > 0\)，则 C \(^{+}\) (Q) 是可分的，且是两个理想的直和，这两个理想都同构于 A 上维数为 \(2^{r-1}\) 的向量空间的自同态代数。

  事实上，由于 Q 是中性的，可以把 E 分解为两个维数为 r 的全奇异子空间 N 与 P 的直和（§ 4 第 2 小节命题 2 的推论 1）。由于 Q 在 N 上的限制为零，C(Q) 的由 N 生成的子代数 S 就等同于 N 的外代数（第 3 小节引理 4）。对 \(n \in \mathbb{N}\)，我们将用 \(e_n'\) 记从 S 到其自身的映射 \(t \to nt\)。

  设 \((n_{1},\ldots,n_{r})\) 是 N 的一个基；我们将用 \((p_{1},\ldots,p_{r})\) 记 P 的使 \(\Phi(n_{i},p_{j})=\delta_{ij}\) 的基（\(\S 4\)，第 2 小节，命题 2）。对 \(p\in P\)，我们将用 \(p'\) 记 N 上的线性型 \(n\to\Phi(n,p)\)，并用 \(i_{p}\) 记 S 的自同态，它是由 T(N) 的自同态 \(i_{p'}\)（与 \(p'\) 相伴，如第 2 小节引理 1 所述）通过过渡到商得到的。由 (5)，我们有

\[
e _ {n} ^ {\prime} \circ i _ {p} + i _ {p} \circ e _ {n} ^ {\prime} = \Phi (n, p) \quad (n \in \mathrm{N}, p \in \mathrm{P}).\tag{12}
\]

  对 \(x = n + p \in \mathbf{E}\)（其中 \(n \in \mathbf{N}\)，\(p \in \mathbf{P}\)），置 \(s(x) = e_n' + i_p\)。显然 \(s\) 是从 \(\mathbf{E}\) 到 \(\mathfrak{L}(\mathbf{S})\) 中的一个线性映射。由于我们有

\[
s (x) ^ {2} = \left(e _ {n} ^ {\prime} + i _ {p}\right) ^ {2} = Q (n) + \Phi (n, p) = Q (x)
\]

  根据 (12) 与引理 1（第 2 小节），s 延拓为 C(Q) 到 ℒ(S) 中的一个同态（我们仍用 s 记它）（第 1 小节命题 1）。我们将证明这个同态是满的；由于 C(Q) 与 ℒ(S) 的维数都是 2²ʳ，这将蕴涵 s 是一个同构，并证明我们的第一个断言。

  事实上，用 I 表示区间 \([1, r]\)。对 I 的每个子集 H，置 \(H' = I - H\)，并用 \(n_{H}\)（相应地，\(p_{H}\)）记诸 \(n_{i}\)（相应地，\(p_{i}\)）（\(i \in H\)，按指标递增顺序排列）的乘积。回忆 \(n_{H}\) 构成 S 的一个基（第 3 小节定理 1）。最后，对 I 的任意两个子集 H、K，置 \(x_{H,K} = n_{H}p_{I}n_{K'}\)。我们将证明元素 \(s(x_{H,K})\)（\(s(C(Q))\) 中的）生成 L(S)。现在，若 \(j \notin H\)，则由引理 1 有 \(s(p_{j})(n_{H}) = i_{p_{j}}(n_{H}) = 0\)，因为诸 \(n_{i}\)（\(i \in H\)）属于 N 上线性型 \(n \to \Phi(n, p_{j})\) 的核；另一方面我们有

\[
s \left(p _ {j}\right) \left(n _ {j} n _ {\mathrm{H}}\right) = \left(i _ {p _ {j}} \circ e _ {n _ {j}} ^ {\prime}\right) \left(n _ {\mathrm{H}}\right) = \Phi \left(p _ {j}, n _ {j}\right) n _ {\mathrm{H}} - n _ {j}. s \left(p _ {j}\right) \left(n _ {\mathrm{H}}\right) = n _ {\mathrm{H}}
\]

（由 (12)）。由于 \(s\) 是一个同态，我们对 I 的任意两个子集 H、K 推出 \(s(p_{\mathrm{K}})(n_{\mathrm{H}}) = 0\)（若 \(\mathrm{K} \not\subset \mathrm{H}\)），以及 \(s(p_{\mathbf{K}}) (n_{\mathbf{H}}) = \pm n_{\mathbf{H} - \mathbf{K}}\)（若 \(\mathrm{K} \subset \mathrm{H}\)）。由于对 \(\mathrm{M} \subset \mathrm{I}\) 与 \(\mathrm{L} \subset \mathrm{I}\)，按定义有 \(s(n_{\mathbf{M}})(n_{\mathbf{L}}) = n_{\mathbf{M}}n_{\mathbf{L}}\)，且 \(n_{\mathbf{M}}n_{\mathbf{L}}\) 在 \(\mathrm{M} \cap \mathrm{L} \neq \emptyset\) 时为零，在相反的情形等于 \(\pm n_{\mathbf{M} \cup \mathbf{L}}\)，我们由前面这些推出，对任意子集 \(\mathrm{H}\)、\(\mathrm{K}\)、\(\mathrm{L}\)（都属于 \(\mathrm{I}\)），\(s(x_{\mathbf{H},\mathbf{K}})(n_{\mathbf{L}}) = s(n_{\mathbf{H}})s(p_{\mathbf{I}})s(n_{\mathbf{K}^{\prime}})(n_{\mathbf{L}})\) 在 \(\mathrm{K} \neq \mathrm{L}\) 时为零，且等于 \(\pm n_{\mathbf{H}}\)（若 \(\mathrm{K} = \mathrm{L}\)）。这表明诸 \(s(x_{\mathbf{H},\mathbf{K}})\) 生成 \(\mathfrak{L}(\mathrm{S})\)，从而完成了第一个断言的证明。

  为证明第二个断言，置 \(S^{+} = S \cap C^{+}\) 与 \(S^{-} = S \cap C^{-}\)；显然 \(S^{+}\)（相应地，\(S^{-}\)）是 S 的由 H 的元素个数为偶数（相应地，奇数）的那些 \(n_{\mathrm{H}}\) 生成的子空间，S 是 \(S^{+}\) 与 \(S^{-}\) 的直和，且 \(s(C^{+})\) 保持 \(S^{+}\) 与 \(S^{-}\) 稳定。因此 \(s\) 把 \(C^{+}\) 映到 \(\mathfrak{L}(S)\) 的一个子代数中，它同构于乘积 \(\mathfrak{L}(S^{+}) \times \mathfrak{L}(S^{-})\)；\(s\) 在 \(C^{+}\) 上的限制是 \(C^{+}\) 到这个子代数上的一个同构，因为 \(s\) 是单射且 \(C^{+}\) 与 \(\mathfrak{L}(S^{+}) \times \mathfrak{L}(S^{-})\) 的维数都是 \(2^{2r-1}\)（第 2 小节定理 1 的推论 1）。证毕。

推论. --- 若 m 为偶数而 Q 的指标任意，则代数 C(Q) 是维数为 \(2^{m}\) 的中心单代数。此外，若 m \textgreater{} 0，子代数 C⁺(Q) 是可分的，且其中心 Z 在 A 上的维数为 2。当 Z 是一个域时，Z 是 A 的一个可分二次扩张且 C⁺(Q) 是单的；否则，Z 是同构于 A 的两个域的直积，而 C⁺(Q) 这时是维数为 \(2^{m-2}\) 的两个单子代数的直积。

  事实上，设 A' 是 A 的代数闭包，Q' 是经纯量扩张从 Q 得到的 \(E' = A' \otimes_{A} E\) 上的二次型。我们已看到 \(C(Q')\) 同构于 \(A' \otimes_{A} C(Q)\)（命题 2），且显然 \(C^{+}(Q')\) 同构于 \(A' \otimes_{A} C^{+}(Q)\)。由于 Q' 是中性的（§ 4 第 2 小节命题 3 的推论 2），本推论是定理 2 与第八章 § 7 的永存性定理的直接推论。

注. --- 1) 由于代数 C(Q) 是单的，它只有一类不可约表示；这些表示称为旋量表示；当注意力固定在这些表示之一（记作 \(\tau\)）上时，\(\tau\) 所在空间的元素称为旋量。若 Q 是中性的，则 \(\tau\) 在 \(C^{+}(Q)\) 上的限制与 \(s\) 的一样，是两个不等价绝对不可约表示之和；承载这两个表示的子空间的元素称为半旋量。在一般情形，若 \(C^{+}(Q)\) 不是单的，则 \(\tau\) 在 \(C^{+}(Q)\) 上的限制由于是忠实的，必须包含属于 \(C^{+}(Q)\) 的两类不可约表示的每一类的子表示，从而是两个不等价绝对不可约表示之和，因为在把纯量扩张到 A 的代数闭包 A' 后情况就是如此。另一方面，若 \(C^{+}(Q)\) 是单的，则它只有一类不可约表示，于是 \(\tau\) 在 \(C^{+}(Q)\) 上的限制是不可约的，因为它在把纯量扩张到 A' 后分解为两个不等价的表示。

\begin{enumerate}
\def\labelenumi{\arabic{enumi})}
\setcounter{enumi}{1}
\tightlist
\item
  假设 A 的特征 \(\neq 2\)，并设 \((x_{1},\ldots,x_{m})\)（\(m=2r\)）是 E 的一个正交基。置
\end{enumerate}

\[
z = 2 ^ {r} x _ {1} \dots x _ {m} \in \mathrm{C(Q)};
\]

  由于 \(x_{i}x_{j} + x_{j}x_{i} = 0\)（对 \(i \neq j\)），我们有 \(zx_{j} = -x_{j}z\)，这蕴涵 \(z\) 属于中心 \(Z\)（\(C^{+}(Q)\) 的）而不属于 \(A\)。我们有

\[
z ^ {2} = 2 ^ {2 r} (- 1) ^ {r} \mathrm{Q} (x _ {1}) \dots \mathrm{Q} (x _ {m}) = (- 1) ^ {r} \mathrm{D}
\]

  其中用 D 表示 \(\Phi\) 关于基 \((x_{j})\) 的判别式（参见习题 9）。

定理 3. --- 假设 E 的维数是奇数 \(m = 2r + 1\)（因而 A 的特征 \(\neq 2\)）。

\begin{enumerate}
\def\labelenumi{\alph{enumi})}
\item
  代数 C \(^{+}\) (Q) 是中心单的。若 Q 有最大指标 r，则 C \(^{+}\) (Q) 同构于 A 上维数为 2 \(^{r}\) 的向量空间的自同态代数。
\item
  代数 C(Q) 是可分的。其中心 Z 的维数为 2，且 C(Q) 同构于 Z⊗ \(_{A}\) C \(^{+}\) (Q)，因而是单的或是两个单子代数的直积。
\end{enumerate}

  事实上，设 \(x_0\) 是 E 的一个非迷向向量，F 是 \(x_0\) 的正交补；用 \(Q_1\) 表示 F 上的二次型 \(y \to -Q(x_0)Q(y)\)；显然 \(Q_1\) 是非退化的。由于 \(x_0y = -yx_0\)（对 \(y \in F\)），我们有 \((x_0y)^2 = -Q(x_0)Q(y) = Q_1(y)\)，因而映射 \(y \to x_0y\)（从 F 到 \(C^+(Q)\) 中）延拓为一个同态 \(h\)（从 \(C(Q_1)\) 到 \(C^+(Q)\) 中）（第 1 小节命题 1）。但 \(C(Q_1)\) 是单的（定理 2）且具有维数 \(2^{2r}\)（与 \(C^{+}(Q)\) 的相同）；由于 \(h(1) = 1\)，这蕴涵 \(h\) 是一个同构。此外，若 \(Q\) 的指标为 \(r\)，则可以选取 \(x_0\) 使 \(Q_1\) 也有指标 \(r\)（\(\S 4\)，第 2 小节，命题 3），这就证明了 \(a\)）。

  现在设 \((x_{1},\ldots ,x_{2r})\) 是 F 的一个正交基；置 \(z = x_0x_1\dots x_{2r}\)。立即可以验证 \(z\) 与 \(x_{j}\)（对 \(j = 0,\dots,2r\)）交换，故属于 C(Q) 的中心。设 Z 是 C(Q) 的由 1 与 \(z\) 生成的子空间；它是 C(Q) 的中心的一个子代数且是 A 的一个二次扩张，因为 \(z\) 是奇的且 \(z^2\) 等于纯量 \((-1)^r\mathrm{Q}(x_0)\ldots \mathrm{Q}(x_{2r})\)。考虑同态 \(\theta\)（从 \(\mathbf{Z}\otimes_{\mathbf{A}}\mathbf{C}^{+}(\mathbf{Q})\) 到 C(Q) 中），它由 \(\theta (u\otimes v) = uv\) 定义。由于 \(z\in \mathbb{C}^{-}\) 且可逆，映射 \(u\to zu\) 是模 \(\mathbf{C}^+\) 到 \(\mathbb{C}^{-}\) 上的一个同构，这蕴涵 \(\theta (\mathrm{Z}\otimes \mathrm{C}^{+})\) 包含 \(\mathbf{C}^{+}\) 与 \(\mathbb{C}^{-}\)，故与 C(Q) 重合。由于 \(\mathbf{Z}\otimes \mathbf{C}^{+}\) 与 C(Q) 有相同的维数 \(2^{2r + 1},\theta\) 是一个同构；这就证明了 \(b)\)，其中用到第八章 § 7 的结果。

注. --- \(\Phi\) 关于基 \((x_{j})_{(j=0,\ldots,2r)}\) 的判别式 D 等于 \(2^{2r+1}Q(x_{0})\ldots Q(x_{2r})\)。因此，Z 由 1 与奇元素 \(z' = 2^{r+1}z\) 生成，其中 \(z'^{2} = (-1)^{r}2D\)。于是代数 C(Q) 是单的当且仅当 \(2(-1)^{r}D\) 在 A 中不是平方。
